\documentclass[12pt, a4paper]{book}
 
%----------------------------------------------------------
%  ENCODAGE & LANGUE
%----------------------------------------------------------
\usepackage[utf8]{inputenc}
\usepackage[T1]{fontenc}
\usepackage[french]{babel}
\frenchbsetup{StandardLists=true}
\usepackage{lmodern}
 
%----------------------------------------------------------
%  MATHS
%----------------------------------------------------------
\usepackage{amsmath, amssymb, mathtools}
 
%----------------------------------------------------------
%  MISE EN PAGE
%----------------------------------------------------------
\usepackage{geometry}
\geometry{a4paper, top=3cm, bottom=3cm, left=3cm, right=3cm, headheight=15pt}
 
%----------------------------------------------------------
%  COULEURS — 3 couleurs cohérentes
%----------------------------------------------------------
\usepackage{xcolor}
 
\definecolor{navy}{RGB}{18, 52, 96}
\definecolor{navylight}{RGB}{235, 241, 250}
\definecolor{forest}{RGB}{26, 80, 50}
\definecolor{forestlight}{RGB}{232, 245, 235}
\definecolor{burgundy}{RGB}{120, 25, 40}
\definecolor{burgundylight}{RGB}{253, 235, 238}
\definecolor{charcoal}{RGB}{52, 58, 64}

\usepackage{pgfplots}
\usepackage{tikz}
\pgfplotsset{compat=1.18}
 
%----------------------------------------------------------
%  BOÎTES  (tcolorbox)
%----------------------------------------------------------
\usepackage[most]{tcolorbox}
 
\newtcolorbox{definition}[1][]{
  breakable,
  before = {\par\vspace{6pt}},
  after  = {\par\vspace{6pt}},
  colback=navylight, colframe=navy,
  boxrule=0pt, leftrule=5pt, arc=0pt,
  left=9pt, right=7pt, top=6pt, bottom=6pt,
  fonttitle=\bfseries\color{navy},
  detach title,
  before upper={\tcbtitle\par\smallskip},
  title={D\'{e}finition \ifx&#1&\else\normalfont\itshape{} : #1\fi}
}
 
\newtcolorbox{theoreme}[1][]{
  breakable,
  before={\par\vspace{6pt}}, after={\par\vspace{6pt}},
  colback=navylight, colframe=navy,
  boxrule=1pt, arc=3pt,
  left=9pt, right=7pt, top=6pt, bottom=6pt,
  coltitle=white, colbacktitle=navy,
  fonttitle=\bfseries,
  title={Th\'{e}or\`{e}me \ifx&#1&\else\normalfont\itshape{} : #1\fi},
  toptitle=3pt, bottomtitle=3pt
}
 
\newtcolorbox{propriete}[1][]{
  breakable,
  before={\par\vspace{6pt}}, after={\par\vspace{6pt}},
  colback=forestlight, colframe=forest,
  boxrule=0pt, leftrule=5pt, arc=0pt,
  left=9pt, right=7pt, top=6pt, bottom=6pt,
  fonttitle=\bfseries\color{forest},
  detach title,
  before upper={\tcbtitle\par\smallskip},
  title={Propri\'{e}t\'{e} \ifx&#1&\else\normalfont\itshape{} : #1\fi}
}
 
\newtcolorbox{methode}[1][]{
  breakable,
  before={\par\vspace{6pt}}, after={\par\vspace{6pt}},
  colback=forestlight, colframe=forest,
  boxrule=1pt, arc=3pt,
  left=9pt, right=7pt, top=6pt, bottom=6pt,
  coltitle=white, colbacktitle=forest,
  fonttitle=\bfseries,
  title={M\'{e}thode \ifx&#1&\else\normalfont\itshape{} : #1\fi},
  toptitle=3pt, bottomtitle=3pt
}
 
\newtcolorbox{rappel}[1][]{
  breakable,
  before={\par\vspace{6pt}}, after={\par\vspace{6pt}},
  colback=burgundylight, colframe=burgundy,
  boxrule=0pt, leftrule=6pt, arc=0pt,
  left=9pt, right=7pt, top=6pt, bottom=6pt,
  fonttitle=\bfseries\color{burgundy},
  detach title,
  before upper={\tcbtitle\par\smallskip},
  title={Rappel \ifx&#1&\else\normalfont\itshape{} : #1\fi}
}
 
\newtcolorbox{attention}{
  breakable,
  before={\par\vspace{6pt}}, after={\par\vspace{6pt}},
  colback=burgundylight, colframe=burgundy,
  boxrule=1pt, arc=3pt,
  left=9pt, right=7pt, top=6pt, bottom=6pt,
  coltitle=white, colbacktitle=burgundy,
  fonttitle=\bfseries,
  title={Attention !},
  toptitle=3pt, bottomtitle=3pt
}
 
\newtcolorbox{exemple}[1][]{
  breakable,
  before={\par\vspace{6pt}}, after={\par\vspace{6pt}},
  colback=gray!7, colframe=charcoal!35,
  boxrule=0pt, leftrule=4pt, arc=0pt,
  left=9pt, right=7pt, top=6pt, bottom=6pt,
  fonttitle=\bfseries\color{charcoal},
  detach title,
  before upper={\tcbtitle\par\smallskip},
  title={Exemple \ifx&#1&\else\normalfont\itshape{} : #1\fi}
}
 
\newtcolorbox{remarque}{
  breakable,
  before={\par},
  colback=white, colframe=gray!40,
  boxrule=0pt, leftrule=3pt, arc=0pt,
  left=9pt, right=7pt, top=4pt, bottom=4pt,
  fontupper=\small\itshape\color{charcoal},
  before upper={\textbf{\small\color{gray!60!black}Remarque.}\enspace}
}
 
\newtcolorbox{exercice}[1][]{
  breakable,
  before={\par\vspace{6pt}}, after={\par\vspace{6pt}},
  colback=white, colframe=navy!45,
  boxrule=0.8pt, arc=3pt,
  left=9pt, right=7pt, top=6pt, bottom=6pt,
  fonttitle=\bfseries\color{navy},
  detach title,
  before upper={\tcbtitle\par\smallskip},
  title={Exercice \ifx&#1&\else #1\fi}
}
 
\newtcolorbox{corollaire}[1][]{
  breakable,
  before={\par\vspace{6pt}}, after={\par\vspace{6pt}},
  colback=forestlight, colframe=forest!60,
  boxrule=0.8pt, arc=3pt,
  left=9pt, right=7pt, top=6pt, bottom=6pt,
  fonttitle=\bfseries\color{forest},
  detach title,
  before upper={\tcbtitle\par\smallskip},
  title={Corollaire \ifx&#1&\else\normalfont\itshape{} : #1\fi}
}
 
%----------------------------------------------------------
%  AUTRES PACKAGES
%----------------------------------------------------------
\usepackage{array, booktabs, colortbl, tabularx}
\usepackage{multirow}       % ← pour le tableau d'en-tête de l'examen
\usepackage{multicol}
\usepackage{lastpage}       % ← pour "Page X / Y" dans l'examen
\usepackage{ifthen}         % ← pour la commande \pts dans l'examen
\usepackage{enumitem}
\renewcommand{\labelitemi}{$\bullet$}
\renewcommand{\labelitemii}{$\circ$}
\renewcommand{\labelitemiii}{$\diamond$}
\usepackage{graphicx}
\usepackage{fancyhdr}
\usepackage{titlesec}
\usepackage{hyperref}
\hypersetup{hidelinks}
 
%----------------------------------------------------------
%  EN-TÊTES & PIEDS DE PAGE
%----------------------------------------------------------
\pagestyle{fancy}
\fancyhf{}
\fancyhead[LE]{\small\color{navy}\textbf{Math\'{e}matiques \textendash{} Bac \'{E}conomie}}
\fancyhead[RO]{\small\color{charcoal}\leftmark}
\fancyfoot[C]{\small\color{charcoal}\thepage}
\renewcommand{\headrulewidth}{0.4pt}
\renewcommand{\headrule}{\color{navy!40}\hrule width\headwidth height\headrulewidth}
 
%----------------------------------------------------------
%  STYLE DES TITRES
%----------------------------------------------------------
\titleformat{\chapter}[display]
  {\normalfont\huge\bfseries\color{navy}}
  {\normalfont\normalsize\color{charcoal!70}\MakeUppercase{\chaptertitlename}\ \thechapter}
  {10pt}{}
  [{\color{navy!30}\titlerule[1.5pt]}]
\titlespacing{\chapter}{0pt}{10pt}{30pt}
 
\titleformat{\section}
  {\normalfont\Large\bfseries\color{navy}}
  {\thesection}{1em}{}
  [{\color{navy!20}\titlerule[0.6pt]}]
\titlespacing{\section}{0pt}{20pt}{10pt}
 
\titleformat{\subsection}
  {\normalfont\large\bfseries\color{forest}}
  {\thesubsection}{1em}{}
\titlespacing{\subsection}{0pt}{14pt}{6pt}
 
\titleformat{\subsubsection}
  {\normalfont\normalsize\bfseries\color{charcoal}}
  {\thesubsubsection}{1em}{}
 
%----------------------------------------------------------
%  COMMANDES MATHÉMATIQUES
%----------------------------------------------------------
\newcommand{\R}{\mathbb{R}}
\newcommand{\N}{\mathbb{N}}
\newcommand{\Z}{\mathbb{Z}}
\newcommand{\Q}{\mathbb{Q}}
\newcommand{\C}{\mathbb{C}}
\newcommand{\dx}{\,dx}
 
%----------------------------------------------------------
%  PAGE DE TITRE — personnalisée sans \maketitle
%----------------------------------------------------------
\begin{document}

\begin{titlepage}
  \centering
  \vspace*{2cm}

  {\color{navy!30}\rule{\linewidth}{1.5pt}}\\[1.2cm]

  {\Huge\bfseries\color{navy} Math\'{e}matiques}\\[0.6cm]

  {\large\color{charcoal}
    2\`{e}me Baccalaur\'{e}at \textemdash{} Sciences \'{E}conomiques}\\[1.2cm]

  {\color{navy!30}\rule{\linewidth}{1.5pt}}

  \vfill

  % Auteur principal
  \begin{center}
    {\large\color{charcoal!60}\textit{Auteur}}\\[0.3cm]
    {\LARGE\bfseries\color{navy} Anass Fawzi}
  \end{center}

  \vspace{1.5cm}

  % Séparateur discret
  {\color{navy!15}\rule{6cm}{0.5pt}}

  \vspace{0.8cm}

  % Co-rédacteur
  \begin{center}
    {\small\color{charcoal!60}\textit{Co-r\'{e}dacteur}}\\[0.3cm]
    {\large\bfseries\color{charcoal} Oussama Sassour}
  \end{center}

  \vspace{1.5cm}

  {\small\color{charcoal!50} Ann\'{e}e scolaire 2024\,--\,2025}

  \vspace{1cm}
\end{titlepage}

%----------------------------------------------------------
%  AVANT-PROPOS
%----------------------------------------------------------
\chapter*{Avant-propos}
\addcontentsline{toc}{chapter}{Avant-propos}

Ce livre de math\'{e}matiques est destin\'{e} aux \'{e}l\`{e}ves de
\textbf{2\`{e}me Baccalaur\'{e}at Sciences \'{E}conomiques} au Maroc.
Il couvre l'ensemble du programme officiel et accompagne l'\'{e}l\`{e}ve
du cours jusqu'aux exercices d'application.

\medskip
Le contenu math\'{e}matique --- d\'{e}finitions, th\'{e}or\`{e}mes,
propri\'{e}t\'{e}s et m\'{e}thodes --- a \'{e}t\'{e} con\c{c}u par
\textbf{Anass Fawzi}.

\medskip
La r\'{e}daction des exemples r\'{e}solus, des exercices d'entra\^{i}nement,
l'organisation en chapitres et la mise en page compl\`{e}te du livre
ont \'{e}t\'{e} assur\'{e}es par \textbf{Oussama Sassour},
co-r\'{e}dacteur de cet ouvrage.

\medskip
Chaque chapitre suit la structure suivante :
\begin{itemize}
  \item \textbf{Rappels} des pr\'{e}requis indispensables
  \item \textbf{D\'{e}finitions} rigoureuses
  \item \textbf{Th\'{e}or\`{e}mes et propri\'{e}t\'{e}s}
  \item \textbf{M\'{e}thodes} et techniques de r\'{e}solution
  \item \textbf{Exemples} r\'{e}solus en d\'{e}tail
  \item \textbf{Exercices} d'entra\^{i}nement
\end{itemize}

\newpage
\tableofcontents
\newpage

%----------------------------------------------------------
%  CHAPITRES
%----------------------------------------------------------
%==========================================================
%  CHAPITRE 1 — LIMITES D'UNE FONCTION
%==========================================================

\chapter{Limites d'une Fonction}

\begin{rappel}
  \begin{itemize}[leftmargin=1.2cm]
    \item Une \textbf{fonction} $f$ associe \`{a} chaque r\'{e}el $x$ de son domaine $D_f$
          un unique r\'{e}el $f(x)$.
    \item Les symboles $+\infty$ et $-\infty$ ne sont \textbf{pas} des r\'{e}els.
    \item On note $\displaystyle\lim_{x \to a} f(x) = \ell$ si $f(x)$
          se rapproche de $\ell$ quand $x$ s'approche de $a$.
  \end{itemize}
\end{rappel}

%----------------------------------------------------------
\section{Limites des fonctions de r\'{e}f\'{e}rence}
%----------------------------------------------------------

\subsection{Limites de \texorpdfstring{$x^n$}{x\^n} et de \texorpdfstring{$\dfrac{1}{x^n}$}{1/x\^n}}

\begin{definition}[Limite en un point et \`{a} l'infini]
  Soit $f$ une fonction et $a \in \R$ ou $a = \pm\infty$.
  \begin{itemize}
    \item $\displaystyle\lim_{x \to a} f(x) = \ell$ signifie que $f(x)$
          peut \^{e}tre rendu aussi proche de $\ell$ que l'on veut,
          pour $x$ suffisamment proche de $a$.
    \item $\displaystyle\lim_{x \to a} f(x) = +\infty$ signifie que $f(x)$
          d\'{e}passe tout r\'{e}el pour $x$ proche de $a$.
  \end{itemize}
\end{definition}

\begin{propriete}[Limites de $x^n$ en $0$ et en $\pm\infty$]
  Soit $n \in \N^*$.

  \bigskip
  \begin{center}
  \renewcommand{\arraystretch}{1.7}
  \begin{tabular}{|c|c|c|c|c|}
    \hline
    \textbf{\color{navy}Fonction}
      & \textbf{\color{navy}$x \to 0$}
      & \textbf{\color{navy}$x \to 0^+$}
      & \textbf{\color{navy}$x \to +\infty$}
      & \textbf{\color{navy}$x \to -\infty$} \\
    \hline\hline
    $x^n$ \; ($n$ pair)   & $0$ & $0$ & $+\infty$ & $+\infty$ \\
    \hline
    $x^n$ \; ($n$ impair) & $0$ & $0$ & $+\infty$ & $-\infty$ \\
    \hline
  \end{tabular}
  \end{center}
\end{propriete}

\begin{propriete}[Limites de $\dfrac{1}{x^n}$ en $0$ et en $\pm\infty$]
  Soit $n \in \N^*$.

  \bigskip
  \begin{center}
  \renewcommand{\arraystretch}{1.9}
  \begin{tabular}{|c|c|c|c|c|c|}
    \hline
    \textbf{\color{navy}Fonction}
      & \textbf{\color{navy}$x \to 0^+$}
      & \textbf{\color{navy}$x \to 0^-$}
      & \textbf{\color{navy}$x \to 0$}
      & \textbf{\color{navy}$x \to +\infty$}
      & \textbf{\color{navy}$x \to -\infty$} \\
    \hline\hline
    $\dfrac{1}{x^n}$ ($n$ pair)
      & $+\infty$ & $+\infty$ & $+\infty$ & $0$ & $0$ \\
    \hline
    $\dfrac{1}{x^n}$ ($n$ impair)
      & $+\infty$ & $-\infty$ & n'existe pas & $0$ & $0$ \\
    \hline
  \end{tabular}
  \end{center}
\end{propriete}

\begin{exemple}[Limites usuelles]
  \[
    \lim_{x \to 0}\, x^3 = 0 \qquad
    \lim_{x \to +\infty} x^4 = +\infty \qquad
    \lim_{x \to -\infty} x^3 = -\infty
  \]
  \[
    \lim_{x \to 0^+} \frac{1}{x} = +\infty \qquad
    \lim_{x \to 0^-} \frac{1}{x} = -\infty \qquad
    \lim_{x \to +\infty} \frac{1}{x^2} = 0
  \]
\end{exemple}

%----------------------------------------------------------
\subsection{Limites des polyn\^{o}mes et fractions rationnelles en \texorpdfstring{$\pm\infty$}{±∞}}
%----------------------------------------------------------

\begin{theoreme}[Limite d'un polyn\^{o}me en $\pm\infty$]
  La limite d'un \textbf{polyn\^{o}me} en $+\infty$ ou $-\infty$ est
  \'{e}gale \`{a} la limite de son \textbf{terme de plus haut degr\'{e}} :
  \[
    \lim_{x \to \pm\infty}
    \bigl(a_n x^n + a_{n-1}x^{n-1} + \cdots + a_0\bigr)
    \;=\;
    \lim_{x \to \pm\infty} a_n\,x^n
  \]
\end{theoreme}

\begin{theoreme}[Limite d'une fraction rationnelle en $\pm\infty$]
  La limite de $\dfrac{P(x)}{Q(x)}$ en $\pm\infty$ est \'{e}gale au
  quotient des \textbf{termes de plus haut degr\'{e}} :
  \[
    \lim_{x \to \pm\infty}
    \frac{a_n x^n + \cdots + a_0}{b_m x^m + \cdots + b_0}
    \;=\;
    \lim_{x \to \pm\infty} \frac{a_n\,x^n}{b_m\,x^m}
  \]
\end{theoreme}

\begin{methode}[Calculer la limite en $\pm\infty$]
  \begin{enumerate}
    \item \textbf{Identifier} le terme de plus haut degr\'{e} au num\'{e}rateur
          et au d\'{e}nominateur.
    \item \textbf{Former} directement le quotient des termes dominants.
    \item \textbf{Conclure} gr\^{a}ce aux limites de r\'{e}f\'{e}rence.
  \end{enumerate}
\end{methode}
\newpage

\begin{exemple}[R\'{e}solus]
  \begin{itemize}
    \item $\displaystyle\lim_{x \to +\infty}(3x^4 - 5x^2 + 2)
          \;=\; \lim_{x \to +\infty} 3x^4 \;=\; +\infty$
    \medskip
    \item $\displaystyle\lim_{x \to -\infty}
          \frac{2x^3 - x + 1}{4x^3 + 7}
          \;=\; \lim_{x \to -\infty} \frac{2x^3}{4x^3}
          \;=\; \frac{1}{2}$
    \medskip
    \item $\displaystyle\lim_{x \to +\infty}
          \frac{x^2 + 1}{x^3 - 2}
          \;=\; \lim_{x \to +\infty} \frac{x^2}{x^3}
          \;=\; \lim_{x \to +\infty} \frac{1}{x} \;=\; 0$
  \end{itemize}
\end{exemple}

%----------------------------------------------------------
\subsection{Limites des fonctions de type \texorpdfstring{$\sqrt{u(x)}$}{sqrt(u(x))}}
%----------------------------------------------------------

\begin{rappel}[Fonction racine carr\'{e}e]
  La fonction $x \mapsto \sqrt{x}$ est d\'{e}finie sur $[0,\,+\infty[$.
  On a $\sqrt{x} \geq 0$ et $\bigl(\sqrt{x}\bigr)^2 = x$ pour tout $x \geq 0$.
\end{rappel}

\begin{propriete}[Limite de $\sqrt{u(x)}$]
  \begin{align*}
    \text{Si } \lim_{x \to x_0} u(x) = \ell \geq 0
    &\;\Longrightarrow\;
    \lim_{x \to x_0} \sqrt{u(x)} = \sqrt{\ell} \\[8pt]
    \text{Si } \lim_{x \to x_0} u(x) = +\infty
    &\;\Longrightarrow\;
    \lim_{x \to x_0} \sqrt{u(x)} = +\infty
  \end{align*}
\end{propriete}

\begin{methode}[Forme ind\'{e}termin\'{e}e $\infty - \infty$ avec une racine]
  On \textbf{multiplie et divise} par l'expression \textbf{conjugu\'{e}e} :
  \[
    \sqrt{a} - \sqrt{b}
    \;=\;
    \frac{\bigl(\sqrt{a}-\sqrt{b}\bigr)\bigl(\sqrt{a}+\sqrt{b}\bigr)}
         {\sqrt{a}+\sqrt{b}}
    \;=\;
    \frac{a - b}{\sqrt{a}+\sqrt{b}}
  \]
\end{methode}

\begin{exemple}[R\'{e}solu : $\infty - \infty$ avec conjugu\'{e}]
  Calculer $\displaystyle\lim_{x \to +\infty}\!\left(\sqrt{x^2+x}-x\right)$.

  \medskip
  \textit{Solution :}
  \[
    \sqrt{x^2+x} - x
    = \frac{x^2+x - x^2}{\sqrt{x^2+x}+x}
    = \frac{x}{\sqrt{x^2+x}+x}
    = \frac{1}{\sqrt{1+\frac{1}{x}}+1}
    \;\xrightarrow[x\to+\infty]{}\; \frac{1}{2}
  \]
\end{exemple}

%----------------------------------------------------------
\newpage
\section{Op\'{e}rations sur les limites}
%----------------------------------------------------------

\begin{definition}[Forme ind\'{e}termin\'{e}e]
  Une \textbf{forme ind\'{e}termin\'{e}e} (F.I.) est une expression dont on
  ne peut pas conclure directement la limite. Les F.I. classiques :
  \[
    \frac{0}{0} \qquad \frac{\infty}{\infty} \qquad
    \infty - \infty \qquad 0 \times \infty
  \]
\end{definition}

\begin{attention}
  Ne jamais \'{e}crire :
  $\;\infty - \infty = 0\;$ ou $\;\dfrac{\infty}{\infty} = 1\;$
  ou $\;0 \times \infty = 0$.
  Ce sont des \textbf{formes ind\'{e}termin\'{e}es}, pas des r\'{e}sultats !
\end{attention}

\begin{theoreme}[Op\'{e}rations sur les limites : Somme]
  \bigskip
  \begin{center}
  \renewcommand{\arraystretch}{1.6}
  \begin{tabular}{|c|c|c|}
    \hline
    \textbf{\color{navy}$\ell = \lim f(x)$}
      & \textbf{\color{navy}$\ell' = \lim g(x)$}
      & \textbf{\color{navy}$\lim(f+g)$} \\
    \hline\hline
    $\ell \in \R$ & $\ell' \in \R$ & $\ell + \ell'$ \\
    \hline
    $\ell \in \R$ & $+\infty$ & $+\infty$ \\
    \hline
    $\ell \in \R$ & $-\infty$ & $-\infty$ \\
    \hline
    $+\infty$ & $+\infty$ & $+\infty$ \\
    \hline
    $-\infty$ & $-\infty$ & $-\infty$ \\
    \hline
    $+\infty$ & $-\infty$ & \textbf{F.I.} \\
    \hline
  \end{tabular}
  \end{center}
\end{theoreme}

\begin{theoreme}[Op\'{e}rations sur les limites : Produit]
  \bigskip
  \begin{center}
  \renewcommand{\arraystretch}{1.6}
  \begin{tabular}{|c|c|c|}
    \hline
    \textbf{\color{navy}$\ell$}
      & \textbf{\color{navy}$\ell'$}
      & \textbf{\color{navy}$\lim(f \times g)$} \\
    \hline\hline
    $\ell \in \R$ & $\ell' \in \R$ & $\ell \times \ell'$ \\
    \hline
    $\ell \in \R^*$ & $\pm\infty$ & $\pm\infty$ (selon les signes) \\
    \hline
    $\pm\infty$ & $\pm\infty$ & $\pm\infty$ (selon les signes) \\
    \hline
    $0$ & $\pm\infty$ & \textbf{F.I.} \\
    \hline
  \end{tabular}
  \end{center}
\end{theoreme}

\begin{theoreme}[Op\'{e}rations sur les limites : Quotient]
  \bigskip
  \begin{center}
  \renewcommand{\arraystretch}{1.9}
  \begin{tabular}{|c|c|c|}
    \hline
    \textbf{\color{navy}$\ell = \lim f(x)$}
      & \textbf{\color{navy}$\ell' = \lim g(x)$}
      & \textbf{\color{navy}$\lim\dfrac{f}{g}$} \\
    \hline\hline
    $\ell \in \R$ & $\ell' \in \R^*$ & $\dfrac{\ell}{\ell'}$ \\
    \hline
    $\ell \in \R$ & $\pm\infty$ & $0$ \\
    \hline
    $\pm\infty$ & $\ell' \in \R^*$ & $\pm\infty$ (selon les signes) \\
    \hline
    $\ell \in \R^*$ ou $\pm\infty$ & $0^+$ ou $0^-$ & $\pm\infty$ (selon les signes) \\
    \hline
    $\pm\infty$ & $\pm\infty$ & \textbf{F.I.} \\
    \hline
    $0$ & $0$ & \textbf{F.I.} \\
    \hline
  \end{tabular}
  \end{center}
\end{theoreme}

%----------------------------------------------------------
\section{Th\'{e}or\`{e}me de comparaison}
%----------------------------------------------------------

\begin{theoreme}[Comparaison : version compl\`{e}te]
  \textbf{(1) Passage \`{a} la limite par in\'{e}galit\'{e} :}
  \[
    f(x) \leq g(x) \;\text{au voisinage de } x_0,\;
    \lim f = \ell,\;\lim g = \ell'
    \quad\Longrightarrow\quad \ell \leq \ell'
  \]
  \textbf{(2) Limite infinie par minoration :}
  \[
    f(x) \leq g(x) \;\text{ et }\;
    \lim_{x \to x_0} f(x) = +\infty
    \quad\Longrightarrow\quad
    \lim_{x \to x_0} g(x) = +\infty
  \]
  \textbf{(3) Limite infinie par majoration :}
  \[
    g(x) \leq f(x) \;\text{ et }\;
    \lim_{x \to x_0} g(x) = -\infty
    \quad\Longrightarrow\quad
    \lim_{x \to x_0} f(x) = -\infty
  \]
\end{theoreme}

\begin{theoreme}[Th\'{e}or\`{e}me des gendarmes]
  Si au voisinage de $x_0$ on a $\;U(x) \leq f(x) \leq V(x)\;$
  et si $\displaystyle\lim_{x \to x_0} U(x) = \lim_{x \to x_0} V(x) = \ell$,
  alors :
  \[
    \lim_{x \to x_0} f(x) = \ell
  \]
\end{theoreme}

\begin{methode}[Appliquer le th\'{e}or\`{e}me des gendarmes]
  \begin{enumerate}
    \item \textbf{Encadrer} $f(x)$ entre deux fonctions $U(x)$ et $V(x)$.
    \item \textbf{V\'{e}rifier} que $U$ et $V$ ont la m\^{e}me limite $\ell$ en $x_0$.
    \item \textbf{Conclure} que $\displaystyle\lim_{x \to x_0} f(x) = \ell$.
  \end{enumerate}
\end{methode}

\begin{exemple}[R\'{e}solu : gendarmes]
  Montrer que $\displaystyle\lim_{x \to +\infty} \frac{\sin x}{x} = 0$.

  \medskip
  Pour tout $x > 0$ : $-1 \leq \sin x \leq 1$, donc
  $\dfrac{-1}{x} \leq \dfrac{\sin x}{x} \leq \dfrac{1}{x}$.
  Comme les deux extr\'{e}mit\'{e}s tendent vers $0$,
  le th\'{e}or\`{e}me des gendarmes donne :
  $\displaystyle\lim_{x \to +\infty} \frac{\sin x}{x} = 0$.
\end{exemple}

%----------------------------------------------------------
\section{R\'{e}capitulatif : lever les formes ind\'{e}termin\'{e}es}
%----------------------------------------------------------

\begin{methode}[Techniques selon la forme ind\'{e}termin\'{e}e]
  \bigskip
  \begin{center}
  \renewcommand{\arraystretch}{1.8}
  \begin{tabular}{|c|l|l|}
    \hline
    \textbf{\color{forest}F.I.}
      & \textbf{\color{forest}Type de fonction}
      & \textbf{\color{forest}Technique \`{a} utiliser} \\
    \hline\hline
    $\dfrac{0}{0}$ ou $\dfrac{\infty}{\infty}$
      & Polyn\^{o}mes / fractions
      & Factoriser par le terme dominant \\
    \hline
    $\dfrac{0}{0}$
      & Avec $\sqrt{\phantom{x}}$
      & Multiplier par l'expression conjugu\'{e}e \\
    \hline
    $\infty - \infty$
      & Avec $\sqrt{\phantom{x}}$
      & Multiplier par le conjugu\'{e} \\
    \hline
    $0 \times \infty$
      & G\'{e}n\'{e}ral
      & R\'{e}\'{e}crire en $\dfrac{0}{1/\infty}$ ou $\dfrac{\infty}{1/0}$ \\
    \hline
  \end{tabular}
  \end{center}
\end{methode}

\begin{exemple}[R\'{e}solu : $\dfrac{0}{0}$ par factorisation]
  Calculer $\displaystyle\lim_{x \to 2} \frac{x^2-4}{x-2}$.

  \medskip
  En $x = 2$ : forme $\dfrac{0}{0}$.
  On factorise : $\dfrac{x^2-4}{x-2} = \dfrac{(x-2)(x+2)}{x-2} = x+2$ \; $(x \neq 2)$.

  Donc $\displaystyle\lim_{x \to 2} \frac{x^2-4}{x-2} = 4$.
\end{exemple}

\begin{exemple}[R\'{e}solu : $\dfrac{\infty}{\infty}$ par terme dominant]
  Calculer $\displaystyle\lim_{x \to +\infty} \frac{3x^2 - x + 5}{x^2 + 4}$.

  \medskip
  On divise par $x^2$ :
  $\dfrac{3x^2-x+5}{x^2+4}
  = \dfrac{3 - \frac{1}{x} + \frac{5}{x^2}}{1 + \frac{4}{x^2}}
  \;\xrightarrow[x\to+\infty]{}\; 3$.
\end{exemple}
 
%----------------------------------------------------------
\section*{Exercices du chapitre}
\addcontentsline{toc}{section}{Exercices : Limites d'une fonction}
%----------------------------------------------------------

\begin{exercice}[1 : Limites de r\'{e}f\'{e}rence]
  Calculer les limites suivantes :
  \[
    \lim_{x \to +\infty} x^5 \qquad
    \lim_{x \to -\infty} x^4 \qquad
    \lim_{x \to 0^+} \frac{1}{x^3} \qquad
    \lim_{x \to 0^-} \frac{1}{x^2}
  \]
\end{exercice}

\begin{exercice}[2 : Polyn\^{o}mes et fractions]
  Calculer les limites en $+\infty$ et $-\infty$ :
  \[
    f(x) = 4x^3 - 2x^2 + 7 \qquad
    g(x) = \frac{2x^2 - 3x + 1}{x^2 + x - 1} \qquad
    h(x) = \frac{x^3-1}{x^2+x}
  \]
\end{exercice}

\begin{exercice}[3 : Formes ind\'{e}termin\'{e}es]
  Lever les F.I. et calculer :
  \begin{multicols}{2}
  \begin{enumerate}
    \item $\displaystyle\lim_{x \to 3} \frac{x^2-9}{x-3}$
    \item $\displaystyle\lim_{x \to +\infty}(\sqrt{x+1}-\sqrt{x})$
    \item $\displaystyle\lim_{x \to +\infty} \frac{x^2-1}{x+1}$
    \item $\displaystyle\lim_{x \to 0} \frac{\sqrt{x+1}-1}{x}$
  \end{enumerate}
  \end{multicols}
\end{exercice}
\begin{exercice}[4 : Th\'{e}or\`{e}me des gendarmes]
  On sait que pour tout $x > 0$ :
  $\dfrac{x^2-1}{x} \leq f(x) \leq \dfrac{x^2+1}{x}$.
  D\'{e}terminer $\displaystyle\lim_{x \to +\infty} \dfrac{f(x)}{x}$.
\end{exercice}
%==========================================================
%  CHAPITRE 2 — DÉRIVABILITÉ
%==========================================================

\chapter{D\'{e}rivabilit\'{e}}

\begin{rappel}[Taux de variation]
  Le \textbf{taux de variation} de $f$ entre $x_0$ et $x$ est :
  \[
    \tau = \frac{f(x) - f(x_0)}{x - x_0}
  \]
  Il repr\'{e}sente le coefficient directeur de la droite passant par
  les points $\bigl(x_0,\,f(x_0)\bigr)$ et $\bigl(x,\,f(x)\bigr)$.
\end{rappel}

%----------------------------------------------------------
\section{D\'{e}rivabilit\'{e} en un point}
%----------------------------------------------------------

\begin{definition}[Fonction d\'{e}rivable en un point]
  On dit que $f$ est \textbf{d\'{e}rivable en $x_0$} si la limite suivante
  existe et est finie :
  \[
    \lim_{x \to x_0} \frac{f(x) - f(x_0)}{x - x_0}
  \]
  Cette limite est appel\'{e}e le \textbf{nombre d\'{e}riv\'{e}} de $f$ en $x_0$,
  not\'{e} $f'(x_0)$.
\end{definition}

\begin{remarque}
  \newline
  On peut aussi \'{e}crire : $f'(x_0) = \displaystyle\lim_{h \to 0}
  \dfrac{f(x_0 + h) - f(x_0)}{h}$ \; en posant $h = x - x_0$.
\end{remarque}

%----------------------------------------------------------
\subsection{D\'{e}rivabilit\'{e} \`{a} droite et \`{a} gauche}
%----------------------------------------------------------

\begin{definition}[D\'{e}rivabilit\'{e} \`{a} droite et \`{a} gauche]
  \begin{itemize}
    \item $f$ est \textbf{d\'{e}rivable \`{a} droite} en $x_0$ si :
    \[
      \lim_{\substack{x \to x_0 \\ x > x_0}}
      \frac{f(x)-f(x_0)}{x-x_0} \text{ existe et est finie}
    \]
    Cette limite est not\'{e}e $f'_d(x_0)$.

    \medskip
    \item $f$ est \textbf{d\'{e}rivable \`{a} gauche} en $x_0$ si :
    \[
      \lim_{\substack{x \to x_0 \\ x < x_0}}
      \frac{f(x)-f(x_0)}{x-x_0} \text{ existe et est finie}
    \]
    Cette limite est not\'{e}e $f'_g(x_0)$.
  \end{itemize}
\end{definition}

\begin{propriete}[Condition de d\'{e}rivabilit\'{e}]
  $f$ est d\'{e}rivable en $x_0$ \textbf{si et seulement si} :
  \[
    f \text{ est d\'{e}rivable \`{a} droite et \`{a} gauche en } x_0
    \quad \text{et} \quad
    f'_d(x_0) = f'_g(x_0)
  \]
\end{propriete}

%----------------------------------------------------------
\section{Tangente et interpr\'{e}tation g\'{e}om\'{e}trique}
%----------------------------------------------------------

\begin{definition}[\'{E}quation de la tangente]
  Soit $f$ d\'{e}rivable en $x_0$. L'\'{e}quation de la \textbf{tangente}
  \`{a} la courbe $\mathcal{C}_f$ au point d'abscisse $x_0$ est :
  \[
    \boxed{y = f'(x_0)(x - x_0) + f(x_0)}
  \]
\end{definition}

\begin{exemple}[R\'{e}solu]
  Soit $f(x) = x^2$. Trouver la tangente en $x_0 = 1$.

  \medskip
  $f'(x) = 2x$ donc $f'(1) = 2$ et $f(1) = 1$.

  L'\'{e}quation de la tangente : $y = 2(x-1) + 1 = 2x - 1$.
\end{exemple}
\newpage
\begin{propriete}[Interpr\'{e}tation g\'{e}om\'{e}trique compl\`{e}te]
  \bigskip
  \begin{center}
  \renewcommand{\arraystretch}{2}
  \begin{tabular}{|p{5cm}|p{8cm}|}
    \hline
    \textbf{\color{navy}Situation}
      & \textbf{\color{navy}Interpr\'{e}tation g\'{e}om\'{e}trique} \\
    \hline\hline
    $f'(x_0) = a \in \R$
      & Tangente au point $A(x_0;\,f(x_0))$, coeff. directeur $a$ \\
    \hline
    $f'(x_0) = 0$
      & Tangente \textbf{horizontale} au point $A$ \\
    \hline
    $f'_d(x_0) = a$
      & Demi-tangente \`{a} droite, coeff. directeur $a$ \\
    \hline
    $f'_d(x_0) = 0$
      & Demi-tangente \textbf{horizontale} \`{a} droite \\
    \hline
    $f'_d(x_0) = +\infty$ ou $-\infty$
      & Demi-tangente \textbf{verticale} \`{a} droite, $f$ \textbf{non d\'{e}rivable} \`{a} droite \\
    \hline
    $f'_g(x_0) = a$
      & Demi-tangente \`{a} gauche, coeff. directeur $a$ \\
    \hline
    $f'_g(x_0) = 0$
      & Demi-tangente \textbf{horizontale} \`{a} gauche \\
    \hline
    $f'_g(x_0) = +\infty$ ou $-\infty$
      & Demi-tangente \textbf{verticale} \`{a} gauche, $f$ \textbf{non d\'{e}rivable} \`{a} gauche \\
    \hline
  \end{tabular}
  \end{center}
\end{propriete}

%----------------------------------------------------------
\section{D\'{e}rivabilit\'{e} et continuit\'{e}}
%----------------------------------------------------------

\begin{theoreme}[D\'{e}rivabilit\'{e} $\Rightarrow$ Continuit\'{e}]
  Si $f$ est d\'{e}rivable en $x_0$, alors $f$ est \textbf{continue} en $x_0$.

  \medskip
  \textbf{Attention :} La r\'{e}ciproque est \textbf{fausse} !
  Une fonction peut \^{e}tre continue en $x_0$ sans y \^{e}tre d\'{e}rivable.
\end{theoreme}

\begin{attention}
  $f$ continue en $x_0$ \textbf{n'implique pas} $f$ d\'{e}rivable en $x_0$.\\[4pt]
  Contre-exemple : $f(x) = |x|$ est continue en $0$ mais \textbf{non d\'{e}rivable} en $0$
  car $f'_d(0) = 1 \neq -1 = f'_g(0)$.
\end{attention}

%----------------------------------------------------------
\newpage
\section{D\'{e}riv\'{e}es des fonctions usuelles}
%----------------------------------------------------------

\begin{propriete}[Tableau des d\'{e}riv\'{e}es usuelles]
  \bigskip
  \begin{center}
  \renewcommand{\arraystretch}{2.2}
  \begin{tabular}{|c|c|c|}
    \hline
    \textbf{\color{forest}Fonction $f(x)$}
      & \textbf{\color{forest}D\'{e}riv\'{e}e $f'(x)$}
      & \textbf{\color{forest}Domaine} \\
    \hline\hline
    $k$ (constante)     & $0$                        & $\R$ \\
    \hline
    $x$                 & $1$                        & $\R$ \\
    \hline
    $x^n$ \; ($n \in \N^*$) & $nx^{n-1}$            & $\R$ \\
    \hline
    $x^r$ \; ($r \in \R$)   & $rx^{r-1}$            & $]0,+\infty[$ \\
    \hline
    $\dfrac{1}{x}$      & $-\dfrac{1}{x^2}$          & $\R^*$ \\
    \hline
    $\sqrt{x}$          & $\dfrac{1}{2\sqrt{x}}$     & $]0,+\infty[$ \\
    \hline
    $e^x$               & $e^x$                      & $\R$ \\
    \hline
    $\ln x$             & $\dfrac{1}{x}$             & $]0,+\infty[$ \\
    \hline
    $\sin x$            & $\cos x$                   & $\R$ \\
    \hline
    $\cos x$            & $-\sin x$                  & $\R$ \\
    \hline
  \end{tabular}
  \end{center}
\end{propriete}

%----------------------------------------------------------
\newpage
\section{Op\'{e}rations sur les d\'{e}riv\'{e}es}
%----------------------------------------------------------

\begin{propriete}[R\`{e}gles de calcul]
  Soient $u$ et $v$ deux fonctions d\'{e}rivables et $k \in \R$.

  \bigskip
  \begin{center}
  \renewcommand{\arraystretch}{2.2}
  \begin{tabular}{|c|c|}
    \hline
    \textbf{\color{navy}Fonction}
      & \textbf{\color{navy}D\'{e}riv\'{e}e} \\
    \hline\hline
    $u + v$              & $u' + v'$ \\
    \hline
    $ku$                 & $ku'$ \\
    \hline
    $u \times v$         & $u'v + uv'$ \\
    \hline
    $\dfrac{u}{v}$       & $\dfrac{u'v - uv'}{v^2}$ \quad ($v \neq 0$) \\
    \hline
    $\dfrac{1}{v}$       & $-\dfrac{v'}{v^2}$ \quad ($v \neq 0$) \\
    \hline
  \end{tabular}
  \end{center}
\end{propriete}

\begin{propriete}[D\'{e}riv\'{e}e d'une fonction compos\'{e}e $u^n$]
  \bigskip
  \begin{center}
  \renewcommand{\arraystretch}{2.2}
  \begin{tabular}{|c|c|}
    \hline
    \textbf{\color{navy}Fonction compos\'{e}e}
      & \textbf{\color{navy}D\'{e}riv\'{e}e} \\
    \hline\hline
    $u^n$ \; ($n \in \Z$)    & $nu'u^{n-1}$ \\
    \hline
    $\sqrt{u}$               & $\dfrac{u'}{2\sqrt{u}}$ \\
    \hline
    $e^u$                    & $u'e^u$ \\
    \hline
    $\ln u$                  & $\dfrac{u'}{u}$ \\
    \hline
    $\sin u$                 & $u'\cos u$ \\
    \hline
    $\cos u$                 & $-u'\sin u$ \\
    \hline
  \end{tabular}
  \end{center}
\end{propriete}

\newpage
\begin{exemple}[R\'{e}solus : op\'{e}rations sur les d\'{e}riv\'{e}es]
  \begin{itemize}
    \item $f(x) = (3x^2 - 1)^4$
    \quad $\Rightarrow$ \quad $u = 3x^2-1$, $u' = 6x$
    \quad $\Rightarrow$ \quad $f'(x) = 4 \cdot 6x \cdot (3x^2-1)^3 = 24x(3x^2-1)^3$

    \medskip
    \item $f(x) = \dfrac{x+1}{x-1}$
    \quad $\Rightarrow$ \quad $u=x+1$, $v=x-1$
    \quad $\Rightarrow$ \quad $f'(x) = \dfrac{1\cdot(x-1) - (x+1)\cdot1}{(x-1)^2} = \dfrac{-2}{(x-1)^2}$

    \medskip
    \item $f(x) = \sqrt{2x+3}$
    \quad $\Rightarrow$ \quad $u = 2x+3$, $u' = 2$
    \quad $\Rightarrow$ \quad $f'(x) = \dfrac{2}{2\sqrt{2x+3}} = \dfrac{1}{\sqrt{2x+3}}$
  \end{itemize}
\end{exemple}

%----------------------------------------------------------
\section{D\'{e}riv\'{e}e et sens de variation}
%----------------------------------------------------------

\begin{theoreme}[Lien entre $f'$ et les variations de $f$]
  Soit $f$ une fonction d\'{e}rivable sur un intervalle $I$.

  \bigskip
  \begin{center}
  \renewcommand{\arraystretch}{1.8}
  \begin{tabular}{|c|c|}
    \hline
    \textbf{\color{navy}Signe de $f'(x)$ sur $I$}
      & \textbf{\color{navy}Variation de $f$ sur $I$} \\
    \hline\hline
    $f'(x) > 0$ pour tout $x \in I$ & $f$ est \textbf{strictement croissante} \\
    \hline
    $f'(x) < 0$ pour tout $x \in I$ & $f$ est \textbf{strictement d\'{e}croissante} \\
    \hline
    $f'(x) = 0$ pour tout $x \in I$ & $f$ est \textbf{constante} \\
    \hline
  \end{tabular}
  \end{center}
\end{theoreme}

\begin{remarque}
  Si $f'(x_0) = 0$ et que $f'$ change de signe en $x_0$,
  alors $f$ admet un \textbf{extremum local} en $x_0$
  (maximum si $f'$ passe de $+$ \`{a} $-$, minimum si $f'$ passe de $-$ \`{a} $+$).
\end{remarque}

\begin{methode}[\'{E}tudier les variations d'une fonction]
  \begin{enumerate}
    \item D\'{e}terminer le \textbf{domaine de d\'{e}finition} $D_f$.
    \item Calculer $f'(x)$.
    \item \'{E}tudier le \textbf{signe de $f'(x)$} sur $D_f$.
    \item Dresser le \textbf{tableau de variations} de $f$.
  \end{enumerate}
\end{methode}

\begin{exemple}[R\'{e}solu : tableau de variations]
  Soit $f(x) = x^3 - 3x + 2$ d\'{e}finie sur $\R$.

  \medskip
  $f'(x) = 3x^2 - 3 = 3(x^2-1) = 3(x-1)(x+1)$

  \medskip
  \begin{center}
  \renewcommand{\arraystretch}{1.8}
  \begin{tabular}{|c|ccccccc|}
    \hline
    $x$ & $-\infty$ && $-1$ && $1$ && $+\infty$ \\
    \hline
    $f'(x)$ && $+$ & $0$ & $-$ & $0$ & $+$ & \\
    \hline
    $f(x)$ & $\nearrow$ && $4$ & $\searrow$ & $0$ & $\nearrow$ & \\
    \hline
  \end{tabular}
  \end{center}

  \medskip
  $f$ admet un \textbf{maximum local} de $4$ en $x = -1$
  et un \textbf{minimum local} de $0$ en $x = 1$.
\end{exemple}

%----------------------------------------------------------
\section*{Exercices du chapitre}
\addcontentsline{toc}{section}{Exercices : D\'{e}rivabilit\'{e}}
%----------------------------------------------------------

\begin{exercice}[1 : Nombre d\'{e}riv\'{e} par la d\'{e}finition]
  En utilisant la d\'{e}finition $f'(x_0) = \displaystyle\lim_{x \to x_0}
  \dfrac{f(x)-f(x_0)}{x-x_0}$, calculer le nombre d\'{e}riv\'{e} de :
  \begin{multicols}{2}
  \begin{enumerate}
    \item $f(x) = x^2 - 3x$ en $x_0 = 2$
    \item $f(x) = \sqrt{x}$ en $x_0 = 4$
    \item $f(x) = \dfrac{1}{x}$ en $x_0 = 1$
    \item $f(x) = x^3$ en $x_0 = -1$
  \end{enumerate}
  \end{multicols}
\end{exercice}

\begin{exercice}[2 : \'{E}quation de la tangente]
  Pour chaque fonction, \'{e}crire l'\'{e}quation de la tangente
  \`{a} la courbe $\mathcal{C}_f$ au point indiqu\'{e} :
  \begin{enumerate}
    \item $f(x) = x^2 + 1$ \quad au point d'abscisse $x_0 = 2$
    \item $f(x) = \sqrt{x}$ \quad au point d'abscisse $x_0 = 9$
    \item $f(x) = e^x$ \quad au point d'abscisse $x_0 = 0$
    \item $f(x) = \ln x$ \quad au point d'abscisse $x_0 = 1$
  \end{enumerate}
\end{exercice}
\newpage
\begin{exercice}[3 : Calcul de d\'{e}riv\'{e}es]
  Calculer la d\'{e}riv\'{e}e des fonctions suivantes :
  \begin{multicols}{2}
  \begin{enumerate}
    \item $f(x) = 3x^4 - 2x^2 + 5$
    \item $f(x) = (2x+1)^5$
    \item $f(x) = \dfrac{x-1}{x+2}$
    \item $f(x) = \sqrt{x^2+1}$
    \item $f(x) = x^2\,e^x$
    \item $f(x) = \ln(3x-1)$
    \item $f(x) = \dfrac{1}{(x+1)^3}$
    \item $f(x) = e^{x^2-1}$
  \end{enumerate}
  \end{multicols}
\end{exercice}

\begin{exercice}[4 : \'{E}tude de variations]
  Pour chacune des fonctions suivantes, calculer $f'(x)$,
  \'{e}tudier son signe et dresser le tableau de variations :
  \begin{enumerate}
    \item $f(x) = 2x^3 - 9x^2 + 12x - 4$ sur $\R$
    \item $f(x) = \dfrac{x^2 - 1}{x^2 + 1}$ sur $\R$
    \item $f(x) = x - 2\sqrt{x}$ sur $[0, +\infty[$
  \end{enumerate}
\end{exercice}

\begin{exercice}[5 : D\'{e}rivabilit\'{e} en un point]
  Soit $f(x) = \begin{cases} x^2 & \text{si } x \leq 1 \\ 2x-1 & \text{si } x > 1 \end{cases}$

  \medskip
  \begin{enumerate}
    \item V\'{e}rifier que $f$ est continue en $x_0 = 1$.
    \item \'{E}tudier la d\'{e}rivabilit\'{e} de $f$ en $x_0 = 1$
          en calculant $f'_d(1)$ et $f'_g(1)$.
    \item Quelle est l'interpr\'{e}tation g\'{e}om\'{e}trique ?
  \end{enumerate}
\end{exercice}

\begin{exercice}[6 : Tangente horizontale]
  Trouver les points de la courbe de $f(x) = x^3 - 6x^2 + 9x + 1$
  en lesquels la tangente est horizontale.
\end{exercice}
%==========================================================
%  CHAPITRE 3 — ÉTUDE DES FONCTIONS
%  Graphes intégrés avec pgfplots + tikz
%==========================================================

\chapter{\'{E}tude des Fonctions}

\begin{rappel}[M\'{e}thode g\'{e}n\'{e}rale d'\'{e}tude d'une fonction]
  L'\'{e}tude compl\`{e}te d'une fonction $f$ comprend :
  \begin{enumerate}
    \item Domaine de d\'{e}finition $D_f$
    \item Parit\'{e}, p\'{e}riodicit\'{e}, sym\'{e}tries
    \item Limites aux bornes et asymptotes
    \item D\'{e}riv\'{e}e $f'$, signe et tableau de variations
    \item Convexit\'{e} : d\'{e}riv\'{e}e seconde $f''$ et points d'inflexion
    \item Trac\'{e} de la courbe repr\'{e}sentative $\mathcal{C}_f$
  \end{enumerate}
\end{rappel}

%----------------------------------------------------------
\section{Sym\'{e}tries d'une courbe}
%----------------------------------------------------------

\subsection{Axe de sym\'{e}trie}

\begin{definition}[Axe de sym\'{e}trie]
  La droite d'\'{e}quation $x = a$ est un \textbf{axe de sym\'{e}trie}
  de la courbe $\mathcal{C}_f$ si et seulement si :
  \[
    \forall x \in D_f \;:\quad (2a - x) \in D_f
    \quad \text{et} \quad
    f(2a - x) = f(x)
  \]
\end{definition}

\begin{remarque}
  Cas particulier : si $a = 0$, la condition devient $f(-x) = f(x)$,
  ce qui signifie que $f$ est une \textbf{fonction paire}
  (sym\'{e}trique par rapport \`{a} l'axe des ordonn\'{e}es).
\end{remarque}

\begin{exemple}[R\'{e}solu : axe de sym\'{e}trie]
  $f(x) = x^2 - 4x + 5$.
  On calcule $f(4 - x) = (4-x)^2 - 4(4-x) + 5 = x^2 - 4x + 5 = f(x)$.

  Donc la droite $x = 2$ est un \textbf{axe de sym\'{e}trie} de $\mathcal{C}_f$.

  \medskip
  \begin{center}
  \begin{tikzpicture}
    \begin{axis}[
      axis lines  = center,
      xlabel      = {$x$},
      ylabel      = {$y$},
      xmin = -0.5, xmax = 4.5,
      ymin = -0.5, ymax = 7,
      xtick       = {0,1,2,3,4},
      ytick       = {0,1,2,3,4,5,6},
      grid        = major,
      grid style  = {gray!20},
      width = 9cm, height = 6cm,
    ]
      % Courbe f(x) = x²-4x+5
      \addplot[navy, very thick, domain=-0.3:4.3, samples=120]
        {x^2 - 4*x + 5};
      \addlegendentry{$f(x)=x^2-4x+5$}

      % Axe de symétrie x=2
      \addplot[burgundy, dashed, thick]
        coordinates {(2,-0.5)(2,7)};
      \addlegendentry{Axe $x=2$}

      % Sommet
      \addplot[mark=*, mark size=3pt, burgundy]
        coordinates {(2,1)};
      \node[right, color=burgundy, font=\small] at (axis cs:2,1) {$(2,\,1)$};
    \end{axis}
  \end{tikzpicture}
  \end{center}
\end{exemple}

%----------------------------------------------------------
\subsection{Centre de sym\'{e}trie}
%----------------------------------------------------------

\begin{definition}[Centre de sym\'{e}trie]
  Le point $I(a,\,b)$ est un \textbf{centre de sym\'{e}trie}
  de la courbe $\mathcal{C}_f$ si et seulement si :
  \[
    \forall x \in D_f \;:\quad (2a - x) \in D_f
    \quad \text{et} \quad
    f(2a - x) + f(x) = 2b
  \]
\end{definition}

\begin{remarque}
  Cas particulier : si $I = (0, 0)$, la condition devient $f(-x) = -f(x)$
  — $f$ est une \textbf{fonction impaire} (sym\'{e}trique par rapport \`{a} l'origine).
\end{remarque}

\begin{methode}[V\'{e}rifier qu'un point est centre de sym\'{e}trie]
  \begin{enumerate}
    \item Calculer directement $f(2a - x) + f(x)$.
    \item Si le r\'{e}sultat vaut $2b$ (constante), alors $I(a,b)$
          est centre de sym\'{e}trie.
  \end{enumerate}
\end{methode}

\begin{exemple}[R\'{e}solu : centre de sym\'{e}trie]
  $f(x) = x^3 - 3x$ sur $\R$.
  V\'{e}rifions que $I(0,\,0)$ est centre de sym\'{e}trie.

  \medskip
  $f(-x) = -x^3 + 3x = -f(x)$,
  donc $f(-x) + f(x) = 0 = 2\times 0$ : $I(0,0)$ est bien centre de sym\'{e}trie.

  \medskip
  \begin{center}
  \begin{tikzpicture}
    \begin{axis}[
      axis lines  = center,
      xlabel      = {$x$},
      ylabel      = {$y$},
      xmin = -2.5, xmax = 2.5,
      ymin = -4,   ymax = 4,
      grid        = major,
      grid style  = {gray!20},
      width = 9cm, height = 7cm,
    ]
      % Courbe f(x) = x³ - 3x
      \addplot[navy, very thick, domain=-2.2:2.2, samples=150]
        {x^3 - 3*x};
      \addlegendentry{$f(x)=x^3-3x$}

      % Centre de symétrie O(0,0)
      \addplot[mark=*, mark size=4pt, burgundy]
        coordinates {(0,0)};
      \node[above right, color=burgundy, font=\small\bfseries]
        at (axis cs:0,0) {$I(0,0)$};

      % Deux points symétriques
      \addplot[mark=square*, mark size=3pt, forest]
        coordinates {(1,-2)(-1,2)};
      \node[right, color=forest, font=\small] at (axis cs:1,-2)  {$(1,-2)$};
      \node[left,  color=forest, font=\small] at (axis cs:-1, 2) {$(-1,2)$};

      % Flèche de symétrie
      \draw[->, forest, dashed] (axis cs:1,-2) -- (axis cs:-1,2);
    \end{axis}
  \end{tikzpicture}
  \end{center}
\end{exemple}

%----------------------------------------------------------
\section{Convexit\'{e} et point d'inflexion}
%----------------------------------------------------------

\begin{definition}[Convexit\'{e} et concavit\'{e}]
  Soit $f$ deux fois d\'{e}rivable sur un intervalle $I$.
  \begin{itemize}
    \item $\mathcal{C}_f$ est \textbf{convexe} sur $I$ si elle est situ\'{e}e
          \textbf{au-dessus} de chacune de ses tangentes :
          \( \forall x \in I,\quad f''(x) \geq 0 \)
    \item $\mathcal{C}_f$ est \textbf{concave} sur $I$ si elle est situ\'{e}e
          \textbf{au-dessous} de chacune de ses tangentes :
          \( \forall x \in I,\quad f''(x) \leq 0 \)
  \end{itemize}
\end{definition}

\begin{definition}[Point d'inflexion]
  Un \textbf{point d'inflexion} est un point o\`{u} $\mathcal{C}_f$
  change de concavit\'{e}.

  \medskip
  \begin{center}
  \renewcommand{\arraystretch}{1.8}
  \begin{tabular}{|p{6.5cm}|p{6.5cm}|}
    \hline
    \textbf{\color{navy}Condition}
      & \textbf{\color{navy}Conclusion} \\
    \hline\hline
    $f''(x_0) = 0$ et $f''$ \textbf{change de signe} en $x_0$
      & $\mathcal{C}_f$ admet un \textbf{point d'inflexion} en $x_0$ \\
    \hline
    $f''(x_0) = 0$ mais \textbf{sans} changement de signe
      & Pas de point d'inflexion en $x_0$ \\
    \hline
  \end{tabular}
  \end{center}
\end{definition}

\begin{attention}
  $f''(x_0) = 0$ \textbf{ne suffit pas} !
  Il faut \textbf{obligatoirement} v\'{e}rifier le changement de signe de $f''$.
\end{attention}

\begin{methode}[\'{E}tudier la convexit\'{e} et trouver les points d'inflexion]
  \begin{enumerate}
    \item Calculer $f''(x)$.
    \item \'{E}tudier le signe de $f''(x)$ sur chaque intervalle.
    \item Les z\'{e}ros o\`{u} $f''$ change de signe $\Rightarrow$ points d'inflexion.
    \item Calculer les coordonn\'{e}es : $\bigl(x_0,\, f(x_0)\bigr)$.
  \end{enumerate}
\end{methode}

\begin{exemple}[R\'{e}solu : convexit\'{e} et point d'inflexion]
  $f(x) = x^3 - 3x^2 + 2$ sur $\R$.

  \medskip
  $f'(x) = 3x^2 - 6x$ \quad et \quad $f''(x) = 6x - 6 = 6(x-1)$.

  \medskip
  \begin{center}
  \renewcommand{\arraystretch}{1.6}
  \begin{tabular}{|c|ccccc|}
    \hline
    $x$           & $-\infty$ && $1$ && $+\infty$ \\
    \hline
    $f''(x)$      && $-$ & $0$ & $+$ & \\
    \hline
    Convexit\'{e} & \multicolumn{2}{c}{Concave $\frown$}
                  && \multicolumn{2}{c|}{Convexe $\smile$} \\
    \hline
  \end{tabular}
  \end{center}

  \medskip
  $f''$ s'annule et change de signe en $x_0 = 1$.
  Le \textbf{point d'inflexion} est $(1,\, f(1)) = (1,\, 0)$.

  \medskip
  \begin{center}
  \begin{tikzpicture}
    \begin{axis}[
      axis lines  = center,
      xlabel      = {$x$},
      ylabel      = {$y$},
      xmin = -0.8, xmax = 3.2,
      ymin = -3,   ymax = 4,
      xtick       = {-0,1,2,3},
      grid        = major,
      grid style  = {gray!20},
      width = 10cm, height = 7cm,
    ]
      % Courbe f(x) = x³ - 3x² + 2
      \addplot[navy, very thick, domain=-0.6:3.1, samples=150]
        {x^3 - 3*x^2 + 2};
      \addlegendentry{$f(x)=x^3-3x^2+2$}

      % Tangente au point d'inflexion (1,0) : f'(1)=-3 → y=-3(x-1)
      \addplot[burgundy, dashed, thick, domain=-0.2:2.5]
        {-3*(x-1)};
      \addlegendentry{Tangente au pt d'inflexion}

      % Point d'inflexion
      \addplot[mark=square*, mark size=4pt, burgundy]
        coordinates {(1,0)};
      \node[above right, color=burgundy, font=\small\bfseries]
        at (axis cs:1,0) {Inflexion $(1,0)$};

      % Annotation concave/convexe
      \node[color=navy!70, font=\small\bfseries] at (axis cs:0.3,2.5)
        {Concave $\frown$};
      \node[color=forest, font=\small\bfseries] at (axis cs:2.5,-1.5)
        {Convexe $\smile$};
    \end{axis}
  \end{tikzpicture}
  \end{center}
\end{exemple}
\newpage
\
%----------------------------------------------------------
\section{Branches infinies et asymptotes}
%----------------------------------------------------------

\begin{definition}[Types d'asymptotes]
  \begin{center}
  \renewcommand{\arraystretch}{2}
  \begin{tabular}{|p{5.5cm}|p{7cm}|}
    \hline
    \textbf{\color{navy}Condition}
      & \textbf{\color{navy}Conclusion} \\
    \hline\hline
    $\displaystyle\lim_{x \to a} f(x) = \pm\infty$
      & Asymptote \textbf{verticale} $x = a$ \\
    \hline
    $\displaystyle\lim_{x \to \pm\infty} f(x) = b \in \R$
      & Asymptote \textbf{horizontale} $y = b$ \\
    \hline
    $\displaystyle\lim_{x \to \pm\infty} \frac{f(x)}{x} = a \neq 0$
    et $\lim [f(x)-ax] = b \in \R$
      & Asymptote \textbf{oblique} $y = ax+b$ \\
    \hline
  \end{tabular}
  \end{center}
\end{definition}

\begin{methode}[Chercher une asymptote oblique en $\pm\infty$]
  \begin{enumerate}
    \item Calculer $a = \displaystyle\lim_{x \to \pm\infty} \dfrac{f(x)}{x}$.
    \item Si $a \in \R^*$, calculer $b = \displaystyle\lim_{x \to \pm\infty}[f(x) - ax]$.
    \item Si $b \in \R$ $\Rightarrow$ asymptote oblique $y = ax+b$.
    \item Si $b = \pm\infty$ $\Rightarrow$ branche parabolique de direction $y=ax$.
  \end{enumerate}
\end{methode}

\begin{exemple}[R\'{e}solu : asymptotes verticale et oblique]
  $f(x) = \dfrac{x^2}{x-1}$ sur $\R\setminus\{1\}$.

  \medskip
  \textbf{A.V. :} $\displaystyle\lim_{x\to 1} f(x) = \pm\infty$
  $\Rightarrow$ asymptote verticale $x = 1$.

  \medskip
  \textbf{A.O. :} $a = \displaystyle\lim_{x\to+\infty}\dfrac{f(x)}{x}
  = \lim\dfrac{x}{x-1} = 1$. \quad
  $b = \lim[f(x)-x] = \lim\dfrac{x^2 - x(x-1)}{x-1} = \lim\dfrac{x}{x-1} = 1$.

  Donc $y = x + 1$ est asymptote oblique.

  \medskip
  \begin{center}
  \begin{tikzpicture}
    \begin{axis}[
      axis lines  = center,
      xlabel      = {$x$},
      ylabel      = {$y$},
      xmin = -3, xmax = 5,
      ymin = -6, ymax = 9,
      grid        = major,
      grid style  = {gray!20},
      width = 10cm, height = 8cm,
      restrict y to domain = -8:10,
    ]
      % Branche gauche
      \addplot[navy, very thick, domain=-2.8:0.75, samples=200]
        {x^2/(x-1)};

      % Branche droite
      \addplot[navy, very thick, domain=1.3:4.8, samples=200]
        {x^2/(x-1)};
      \addlegendentry{$f(x)=\dfrac{x^2}{x-1}$}

      % Asymptote verticale x=1
      \addplot[burgundy, dashed, thick]
        coordinates {(1,-6)(1,9)};
      \addlegendentry{A.V. : $x=1$}

      % Asymptote oblique y = x+1
      \addplot[forest, dashed, thick, domain=-3:5]
        {x+1};
      \addlegendentry{A.O. : $y=x+1$}
    \end{axis}
  \end{tikzpicture}
  \end{center}
\end{exemple}

%----------------------------------------------------------
\section{Fonction r\'{e}ciproque}
%----------------------------------------------------------

\begin{definition}[Fonction r\'{e}ciproque]
  Si $f$ est \textbf{continue et strictement monotone} sur $I$,
  alors $f$ admet une \textbf{fonction r\'{e}ciproque} $f^{-1}$
  d\'{e}finie sur $f(I)$, v\'{e}rifiant :
  \[
    f^{-1}(x) = y \;\iff\; f(y) = x
  \]
\end{definition}

\begin{propriete}[Propri\'{e}t\'{e}s de $f^{-1}$]
  \begin{itemize}
    \item $f^{-1}$ est \textbf{continue} sur $f(I)$.
    \item $f^{-1}$ varie dans le \textbf{m\^{e}me sens} que $f$.
    \item Si $f'$ ne s'annule pas sur $I$ :
    \[
      \bigl(f^{-1}\bigr)'(x) = \frac{1}{f'\bigl(f^{-1}(x)\bigr)}
      \qquad \text{et} \qquad
      \bigl(f^{-1}\bigr)'(y_0) = \frac{1}{f'(x_0)}
      \quad \text{avec } y_0 = f(x_0)
    \]
  \end{itemize}
\end{propriete}

\begin{methode}[Trouver l'expression de $f^{-1}(x)$]
  \begin{enumerate}
    \item Poser $f(y) = x$.
    \item \textbf{R\'{e}soudre} cette \'{e}quation en $y$.
    \item La solution $y$ est $f^{-1}(x)$.
  \end{enumerate}
\end{methode}

\begin{exemple}[R\'{e}solu : $f$ et $f^{-1}$ sym\'{e}triques par rapport \`{a} $y=x$]
  $f(x) = e^x$ sur $\R$, donc $f^{-1}(x) = \ln x$ sur $]0,+\infty[$.

  \medskip
  \begin{center}
  \begin{tikzpicture}
    \begin{axis}[
      axis lines  = center,
      xlabel      = {$x$},
      ylabel      = {$y$},
      xmin = -2.5, xmax = 4,
      ymin = -2.5, ymax = 4,
      grid        = major,
      grid style  = {gray!20},
      width = 9cm, height = 9cm,
    ]
      % f(x) = e^x
      \addplot[navy, very thick, domain=-2.4:1.38, samples=120]
        {exp(x)};
      \addlegendentry{$f(x)=e^x$}

      % f⁻¹(x) = ln(x)
      \addplot[forest, very thick, domain=0.09:4, samples=120]
        {ln(x)};
      \addlegendentry{$f^{-1}(x)=\ln x$}

      % Première bissectrice y = x
      \addplot[burgundy, dashed, thick, domain=-2:4]
        {x};
      \addlegendentry{Bissectrice $y=x$}

      % Point commun (0,1) et (1,0)
      \addplot[mark=*, mark size=3pt, navy]
        coordinates {(0,1)};
      \node[left, font=\small, navy] at (axis cs:0,1) {$(0,1)$};

      \addplot[mark=*, mark size=3pt, forest]
        coordinates {(1,0)};
      \node[below, font=\small, forest] at (axis cs:1,0) {$(1,0)$};
    \end{axis}
  \end{tikzpicture}
  \end{center}
\end{exemple}

\begin{propriete}[Correspondance des asymptotes entre $f$ et $f^{-1}$]
  \begin{center}
  \renewcommand{\arraystretch}{1.8}
  \begin{tabular}{|p{5.5cm}|p{5.5cm}|}
    \hline
    \textbf{\color{navy}$\mathcal{C}_f$ admet\ldots}
      & \textbf{\color{navy}$\mathcal{C}_{f^{-1}}$ admet\ldots} \\
    \hline\hline
    A. verticale $x = a$     & A. horizontale $y = a$ \\
    \hline
    A. horizontale $y = a$   & A. verticale $x = a$ \\
    \hline
    A. oblique $y = ax + b$  & A. oblique $y = \dfrac{1}{a}x - \dfrac{b}{a}$ \\
    \hline
    Tangente verticale       & Tangente horizontale \\
    \hline
    Tangente horizontale     & Tangente verticale \\
    \hline
  \end{tabular}
  \end{center}
\end{propriete}

\newpage
%----------------------------------------------------------
\section*{Exercices du chapitre}
\addcontentsline{toc}{section}{Exercices : \'{E}tude des fonctions}
%----------------------------------------------------------

\begin{exercice}[1 : Sym\'{e}tries]
  Pour chaque fonction, d\'{e}terminer si la courbe poss\`{e}de
  un axe ou un centre de sym\'{e}trie :
  \begin{enumerate}
    \item $f(x) = x^4 - 2x^2 + 1$
    \item $f(x) = x^3 - 3x$
    \item $f(x) = \dfrac{1}{x-2} + 3$
  \end{enumerate}
\end{exercice}

\begin{exercice}[2 : Convexit\'{e} et points d'inflexion]
  \'{E}tudier la convexit\'{e} et d\'{e}terminer les \'{e}ventuels points d'inflexion :
  \begin{enumerate}
    \item $f(x) = x^3 - 6x^2 + 9x$ sur $\R$
    \item $f(x) = e^x - x$ sur $\R$
    \item $f(x) = \ln x$ sur $]0,+\infty[$
  \end{enumerate}
\end{exercice}

\begin{exercice}[3 : Asymptotes]
  \'{E}tudier les branches infinies de :
  \begin{enumerate}
    \item $f(x) = \dfrac{x^2-1}{x+1}$ \quad (asymptote oblique ?)
    \item $f(x) = \dfrac{2x}{x-3}$ \quad (asymptotes verticale et horizontale)
    \item $f(x) = x + \sqrt{x^2+1}$ \quad (asymptote oblique en $\pm\infty$)
  \end{enumerate}
\end{exercice}

\begin{exercice}[4 : Fonction r\'{e}ciproque]
  D\'{e}terminer $f^{-1}$ et son domaine :
  \begin{enumerate}
    \item $f(x) = 3x - 5$ sur $\R$
    \item $f(x) = x^2 + 1$ sur $[0,+\infty[$
    \item $f(x) = e^{x+1}$ sur $\R$
    \item $f(x) = \dfrac{x+1}{x-1}$ sur $]1,+\infty[$
  \end{enumerate}
\end{exercice}

\begin{exercice}[5 : \'{E}tude compl\`{e}te avec trac\'{e}]
  Soit $f(x) = \dfrac{x^2}{x-1}$ sur $\R\setminus\{1\}$.
  \begin{enumerate}
    \item Limites aux bornes et asymptotes.
    \item $f'(x)$, signe et tableau de variations.
    \item $f''(x)$, convexit\'{e} et point d'inflexion.
    \item Tracer $\mathcal{C}_f$ avec ses asymptotes.
  \end{enumerate}
\end{exercice}
%==========================================================
%  CHAPITRE 4 — SUITES NUMÉRIQUES
%==========================================================

\chapter{Suites Num\'{e}riques}

\begin{rappel}[Vocabulaire de base]
  \begin{itemize}
    \item Une \textbf{suite num\'{e}rique} est une fonction $u : \N \to \R$,
          not\'{e}e $(u_n)_{n \in \N}$.
    \item On peut d\'{e}finir une suite par une \textbf{formule explicite}
          $u_n = f(n)$ ou par une \textbf{relation de r\'{e}currence}
          $u_{n+1} = f(u_n)$.
    \item Le \textbf{terme g\'{e}n\'{e}ral} $u_n$ est la valeur de la suite au rang $n$.
  \end{itemize}
\end{rappel}

%----------------------------------------------------------
\section{Suites arithm\'{e}tiques et g\'{e}om\'{e}triques}
%----------------------------------------------------------

\begin{definition}[Suite arithm\'{e}tique]
  Une suite $(u_n)$ est \textbf{arithm\'{e}tique} de raison $r$ si :
  \[
    \forall n \in \N, \quad u_{n+1} = u_n + r
  \]
\end{definition}

\begin{definition}[Suite g\'{e}om\'{e}trique]
  Une suite $(u_n)$ est \textbf{g\'{e}om\'{e}trique} de raison $q$ si :
  \[
    \forall n \in \N, \quad u_{n+1} = q \cdot u_n \qquad (q \neq 0,\ u_0 \neq 0)
  \]
\end{definition}
\newpage 
\begin{propriete}[Tableau comparatif]
  \bigskip
  \renewcommand{\arraystretch}{2.2}
  \begin{tabularx}{\linewidth}{|>{\bfseries}p{3.2cm}|>{\centering\arraybackslash}X|>{\centering\arraybackslash}X|}
    \hline
    \rowcolor{white}
      & \textbf{\color{navy}Suite arithm\'{e}tique}
      & \textbf{\color{forest}Suite g\'{e}om\'{e}trique} \\
    \hline\hline
    D\'{e}finition
      & $u_{n+1} = u_n + r$
      & $u_{n+1} = q \cdot u_n$ \\
    \hline
    Terme g\'{e}n\'{e}ral
      & $u_n = u_p + (n-p)\,r$
      & $u_n = u_p \cdot q^{n-p}$ \\
    \hline
    Somme de termes cons\'{e}cutifs
      & $S = \dfrac{(n-p+1)(u_p + u_n)}{2}$
      & $S = u_p \cdot \dfrac{1 - q^{n-p+1}}{1-q}$ \; ($q \neq 1$) \\
    \hline
    $a$, $b$, $c$ cons\'{e}cutifs
      & $2b = a + c$
      & $b^2 = a \cdot c$ \\
    \hline
  \end{tabularx}
\end{propriete}

\begin{remarque}
  Pour la somme g\'{e}om\'{e}trique, si $q = 1$ alors $S = (n-p+1)\,u_p$.

  \medskip
  Formule utile : $1 + q + q^2 + \cdots + q^n = \dfrac{1 - q^{n+1}}{1-q}$
  pour $q \neq 1$.
\end{remarque}

\begin{exemple}[R\'{e}solu : suite arithm\'{e}tique]
  $(u_n)$ arithm\'{e}tique avec $u_0 = 3$ et $r = 5$.

  \medskip
  \textbf{Terme g\'{e}n\'{e}ral :} $u_n = 3 + 5n$.

  \medskip
  \textbf{Somme :} $u_0 + u_1 + \cdots + u_{10}
  = \dfrac{11 \times (u_0 + u_{10})}{2}
  = \dfrac{11 \times (3 + 53)}{2} = \dfrac{11 \times 56}{2} = 308$.
\end{exemple}

\begin{exemple}[R\'{e}solu : suite g\'{e}om\'{e}trique]
  $(u_n)$ g\'{e}om\'{e}trique avec $u_0 = 2$ et $q = 3$.

  \medskip
  \textbf{Terme g\'{e}n\'{e}ral :} $u_n = 2 \times 3^n$.

  \medskip
  \textbf{Somme :} $u_0 + u_1 + \cdots + u_4
  = 2 \times \dfrac{1 - 3^5}{1 - 3}
  = 2 \times \dfrac{1 - 243}{-2} = 242$.
\end{exemple}
\newpage
%----------------------------------------------------------
\section{Monotonie et bornitude}
%----------------------------------------------------------

\begin{definition}[Monotonie d'une suite]
  Soit $(u_n)$ une suite num\'{e}rique.
  \begin{itemize}
    \item $(u_n)$ est \textbf{croissante} si :
          $\forall n \in \N,\quad u_{n+1} \geq u_n$
    \item $(u_n)$ est \textbf{d\'{e}croissante} si :
          $\forall n \in \N,\quad u_{n+1} \leq u_n$
    \item $(u_n)$ est \textbf{constante} si :
          $\forall n \in \N,\quad u_{n+1} = u_n$
  \end{itemize}
\end{definition}

\begin{methode}[\'{E}tudier la monotonie d'une suite]
  On peut \'{e}tudier le signe de $u_{n+1} - u_n$, ou,
  si les termes sont positifs, comparer $\dfrac{u_{n+1}}{u_n}$ \`{a} $1$.

  \medskip
  \begin{center}
  \renewcommand{\arraystretch}{1.8}
  \begin{tabular}{|c|c|}
    \hline
    \textbf{\color{navy}Condition}
      & \textbf{\color{navy}Conclusion} \\
    \hline\hline
    $u_{n+1} - u_n > 0$
      & $(u_n)$ \textbf{croissante} \\
    \hline
    $u_{n+1} - u_n < 0$
      & $(u_n)$ \textbf{d\'{e}croissante} \\
    \hline
    $\dfrac{u_{n+1}}{u_n} > 1$ \; ($u_n > 0$)
      & $(u_n)$ \textbf{croissante} \\
    \hline
    $\dfrac{u_{n+1}}{u_n} < 1$ \; ($u_n > 0$)
      & $(u_n)$ \textbf{d\'{e}croissante} \\
    \hline
  \end{tabular}
  \end{center}
\end{methode}

\begin{definition}[Suite born\'{e}e]
  Soit $(u_n)$ une suite num\'{e}rique.
  \begin{itemize}
    \item $(u_n)$ est \textbf{major\'{e}e} par $M$ si :
          $\forall n \in \N,\quad u_n \leq M$
    \item $(u_n)$ est \textbf{minor\'{e}e} par $m$ si :
          $\forall n \in \N,\quad u_n \geq m$
    \item $(u_n)$ est \textbf{born\'{e}e} si elle est \`{a} la fois
          major\'{e}e et minor\'{e}e.
  \end{itemize}
\end{definition}

\begin{exemple}[R\'{e}solu : monotonie et bornitude]
  $u_n = \dfrac{n}{n+1}$.

  \medskip
  $u_{n+1} - u_n
  = \dfrac{n+1}{n+2} - \dfrac{n}{n+1}
  = \dfrac{(n+1)^2 - n(n+2)}{(n+2)(n+1)}
  = \dfrac{1}{(n+2)(n+1)} > 0$.

  Donc $(u_n)$ est \textbf{croissante}.

  De plus $0 \leq u_n < 1$ pour tout $n$ :
  $(u_n)$ est \textbf{born\'{e}e} (minor\'{e}e par $0$, major\'{e}e par $1$).
\end{exemple}

%----------------------------------------------------------
\section{Limite d'une suite}
%----------------------------------------------------------

\begin{definition}[Convergence et divergence]
  \begin{itemize}
    \item $(u_n)$ \textbf{converge} vers $\ell \in \R$ si $u_n$
          peut \^{e}tre rendu aussi proche de $\ell$ que l'on veut
          pour $n$ suffisamment grand. On \'{e}crit :
          $\displaystyle\lim_{n \to +\infty} u_n = \ell$
    \item $(u_n)$ \textbf{diverge vers $+\infty$} si $u_n$ d\'{e}passe
          tout r\'{e}el pour $n$ grand.
  \end{itemize}
\end{definition}

\begin{propriete}[Limite de $n^\alpha$ avec $\alpha \in \R^*$]
  \[
    \lim_{n \to +\infty} n^\alpha = +\infty \quad (\alpha > 0)
    \qquad\qquad
    \lim_{n \to +\infty} n^\alpha = 0 \quad (\alpha < 0)
  \]
\end{propriete}

\begin{propriete}[Limite de $q^n$ selon la valeur de $q$]
  \bigskip
  \begin{center}
  \renewcommand{\arraystretch}{1.8}
  \begin{tabular}{|c|c|c|c|c|}
    \hline
    \textbf{\color{navy}$q < -1$}
      & \textbf{\color{navy}$q = -1$}
      & \textbf{\color{navy}$-1 < q < 1$}
      & \textbf{\color{navy}$q = 1$}
      & \textbf{\color{navy}$q > 1$} \\
    \hline\hline
    Pas de limite & Pas de limite & $0$ & $1$ & $+\infty$ \\
    \hline
  \end{tabular}
  \end{center}
\end{propriete}

\begin{propriete}[Th\'{e}or\`{e}me des gendarmes pour les suites]
  Si $\forall n \in \N : \; U_n \leq u_n \leq V_n$
  et si $\displaystyle\lim_{n\to+\infty} U_n
       = \lim_{n\to+\infty} V_n = \ell$,
  alors :
  \[
    \lim_{n \to +\infty} u_n = \ell
  \]
\end{propriete}

\begin{propriete}[Th\'{e}or\`{e}me de comparaison]
  \begin{itemize}
    \item Si $u_n \leq v_n$ et $\displaystyle\lim u_n = +\infty$,
          alors $\displaystyle\lim v_n = +\infty$.
    \item Si $u_n \leq v_n$ pour tout $n$ et $(u_n)$, $(v_n)$
          convergent, alors $\displaystyle\lim u_n \leq \lim v_n$.
  \end{itemize}
\end{propriete}
\newpage

%----------------------------------------------------------
\section{Crit\`{e}res de convergence}
%----------------------------------------------------------

\begin{theoreme}[Suites monotones et born\'{e}es]
  \begin{itemize}
    \item Toute suite \textbf{croissante et major\'{e}e} est \textbf{convergente}.
    \item Toute suite \textbf{d\'{e}croissante et minor\'{e}e} est \textbf{convergente}.
  \end{itemize}
\end{theoreme}

\begin{attention}
  Ces th\'{e}or\`{e}mes affirment l'existence de la limite,
  mais \textbf{ne donnent pas sa valeur}.
  Pour trouver $\ell$, on utilise les m\'{e}thodes des sections suivantes.
\end{attention}

%----------------------------------------------------------
\section{Suite de type \texorpdfstring{$v_n = f(u_n)$}{v\_n = f(u\_n)}}
%----------------------------------------------------------

\begin{theoreme}[Suite compos\'{e}e]
  Si $(u_n)$ est convergente de limite $\ell$ et si $f$ est
  \textbf{continue en $\ell$}, alors la suite $(v_n)$ d\'{e}finie par
  $v_n = f(u_n)$ est convergente et :
  \[
    \lim_{n \to +\infty} v_n = f(\ell)
  \]
\end{theoreme}

\begin{exemple}[R\'{e}solu]
  $u_n = \dfrac{1}{n}$ (converge vers $0$).
  Soit $v_n = e^{u_n} = e^{1/n}$.

  $f(x) = e^x$ est continue en $0$, donc :
  $\displaystyle\lim_{n\to+\infty} v_n = e^0 = 1$.
\end{exemple}

%----------------------------------------------------------
\section{Suite de type \texorpdfstring{$u_{n+1} = f(u_n)$}{u\_{n+1} = f(u\_n)}}
%----------------------------------------------------------

\begin{theoreme}[Suite r\'{e}currente]
  Soit $(u_n)$ d\'{e}finie par $u_0 = a$ et $u_{n+1} = f(u_n)$.
  Si :
  \begin{itemize}
    \item $f$ est \textbf{continue} sur un intervalle $I$
    \item $f(I) \subset I$ \quad (stabilit\'{e})
    \item $a \in I$
    \item $(u_n)$ est \textbf{convergente} de limite $\ell$
  \end{itemize}
  alors $\ell$ est \textbf{solution de l'\'{e}quation} $f(x) = x$
  (point fixe de $f$).
\end{theoreme}

\begin{methode}[\'{E}tudier une suite r\'{e}currente $u_{n+1} = f(u_n)$]
  \begin{enumerate}
    \item \textbf{Montrer} que $(u_n)$ est dans un intervalle $I$
          stable par $f$ (par r\'{e}currence).
    \item \textbf{Montrer} la monotonie (croissante ou d\'{e}croissante).
    \item \textbf{Conclure} \`{a} la convergence (born\'{e}e + monotone).
    \item \textbf{Trouver} la limite $\ell$ en r\'{e}solvant $f(\ell) = \ell$.
  \end{enumerate}
\end{methode}

\begin{exemple}[R\'{e}solu : suite r\'{e}currente]
  $u_0 = 0$ et $u_{n+1} = \dfrac{u_n + 2}{2}$.

  \medskip
  \textbf{\'{E}tape 1 :} Montrons par r\'{e}currence que $0 \leq u_n \leq 2$.

  $u_0 = 0 \in [0,2]$. Si $u_n \in [0,2]$ alors
  $u_{n+1} = \dfrac{u_n+2}{2} \in [1,2] \subset [0,2]$. $\checkmark$

  \medskip
  \textbf{\'{E}tape 2 :}
  $u_{n+1} - u_n = \dfrac{u_n+2}{2} - u_n = \dfrac{2 - u_n}{2} \geq 0$
  car $u_n \leq 2$.
  Donc $(u_n)$ est \textbf{croissante}.

  \medskip
  \textbf{\'{E}tape 3 :} Croissante et major\'{e}e par $2$ $\Rightarrow$
  $(u_n)$ \textbf{converge}.

  \medskip
  \textbf{\'{E}tape 4 :} $\ell = f(\ell)$ :
  $\ell = \dfrac{\ell + 2}{2} \Rightarrow 2\ell = \ell + 2 \Rightarrow \ell = 2$.
\end{exemple}
\newpage
%----------------------------------------------------------
\section{Repr\'{e}sentation graphique d'une suite r\'{e}currente}
%----------------------------------------------------------
\begin{propriete}[Toile d'araign\'{e}e]
  La suite $u_{n+1} = f(u_n)$ se visualise graphiquement par
  la \textbf{m\'{e}thode de la toile d'araign\'{e}e} :
  \begin{enumerate}
    \item Tracer $\mathcal{C}_f$ (courbe de $f$) et la droite $y = x$.
    \item Partir du point $(u_0, 0)$ sur l'axe des abscisses.
    \item Monter verticalement jusqu'\`{a} $\mathcal{C}_f$ :
          on obtient $(u_0, f(u_0)) = (u_0, u_1)$.
    \item Aller horizontalement jusqu'\`{a} la droite $y=x$ :
          on obtient $(u_1, u_1)$.
    \item R\'{e}p\'{e}ter les \'{e}tapes 3--4.
  \end{enumerate}
\end{propriete}

\begin{exemple}[Toile d'araign\'{e}e : $u_{n+1} = \frac{u_n+2}{2}$, $u_0 = 0$]
  \begin{center}
  \begin{tikzpicture}
    \begin{axis}[
      axis lines  = center,
      xlabel      = {$x$},
      ylabel      = {$y$},
      xmin = -0.2, xmax = 2.8,
      ymin = -0.2, ymax = 2.8,
      xtick       = {0,0.5,1,1.5,2,2.5},
      ytick       = {0,0.5,1,1.5,2,2.5},
      grid        = major,
      grid style  = {gray!20},
      width = 10cm, height = 10cm,
    ]
      % Courbe f(x) = (x+2)/2
      \addplot[navy, very thick, domain=0:2.8]
        {(x+2)/2};
      \addlegendentry{$f(x)=\frac{x+2}{2}$}

      % Droite y = x
      \addplot[burgundy, dashed, thick, domain=0:2.8]
        {x};
      \addlegendentry{$y = x$}

      % Point fixe ℓ = 2
      \addplot[mark=*, mark size=4pt, burgundy]
        coordinates {(2,2)};
      \node[above left, burgundy, font=\small\bfseries]
        at (axis cs:2,2) {$\ell=2$};

      % Toile d'araignée : u0=0
      % (0,0) → (0, f(0)=1)
      \addplot[forest, thick]
        coordinates {(0,0)(0,1)};
      % (0,1) → (1,1)
      \addplot[forest, thick]
        coordinates {(0,1)(1,1)};
      % (1,1) → (1, f(1)=1.5)
      \addplot[forest, thick]
        coordinates {(1,1)(1,1.5)};
      % (1,1.5) → (1.5,1.5)
      \addplot[forest, thick]
        coordinates {(1,1.5)(1.5,1.5)};
      % (1.5,1.5) → (1.5, f(1.5)=1.75)
      \addplot[forest, thick]
        coordinates {(1.5,1.5)(1.5,1.75)};
      % (1.5,1.75) → (1.75,1.75)
      \addplot[forest, thick]
        coordinates {(1.5,1.75)(1.75,1.75)};
      % (1.75,1.75) → (1.75, f(1.75)=1.875)
      \addplot[forest, thick]
        coordinates {(1.75,1.75)(1.75,1.875)};
      % (1.75,1.875) → (1.875,1.875)
      \addplot[forest, thick]
        coordinates {(1.75,1.875)(1.875,1.875)};

      % Points u0, u1, u2, u3
      \addplot[mark=|, mark size=5pt, forest]
        coordinates {(0,0)(1,0)(1.5,0)(1.75,0)};
      \node[below, forest, font=\footnotesize] at (axis cs:0,0)    {$u_0$};
      \node[below, forest, font=\footnotesize] at (axis cs:1,0)    {$u_1$};
      \node[below, forest, font=\footnotesize] at (axis cs:1.5,0)  {$u_2$};
      \node[below, forest, font=\footnotesize] at (axis cs:1.75,0) {$u_3$};
    \end{axis}
  \end{tikzpicture}
  \end{center}
\end{exemple}

\newpage
%----------------------------------------------------------
\section*{Exercices du chapitre}
\addcontentsline{toc}{section}{Exercices : Suites num\'{e}riques}
%----------------------------------------------------------

\begin{exercice}[1 : Suite arithm\'{e}tique]
  $(u_n)$ est arithm\'{e}tique avec $u_2 = 7$ et $u_5 = 16$.
  \begin{enumerate}
    \item D\'{e}terminer la raison $r$ et le premier terme $u_0$.
    \item Exprimer $u_n$ en fonction de $n$.
    \item Calculer $S = u_0 + u_1 + \cdots + u_{20}$.
  \end{enumerate}
\end{exercice}

\begin{exercice}[2 : Suite g\'{e}om\'{e}trique]
  $(u_n)$ est g\'{e}om\'{e}trique avec $u_0 = 5$ et $u_3 = 40$.
  \begin{enumerate}
    \item D\'{e}terminer la raison $q$.
    \item Exprimer $u_n$ en fonction de $n$.
    \item Calculer $S = u_0 + u_1 + \cdots + u_7$.
  \end{enumerate}
\end{exercice}

\begin{exercice}[3 : Monotonie et bornitude]
  Pour chaque suite, \'{e}tudier la monotonie et la bornitude :
  \begin{enumerate}
    \item $u_n = \dfrac{2n+1}{n+2}$
    \item $u_n = 2^n - n$
    \item $u_n = \dfrac{(-1)^n}{n+1}$
  \end{enumerate}
\end{exercice}

\begin{exercice}[4 : Limite d'une suite]
  Calculer les limites suivantes :
  \begin{multicols}{2}
  \begin{enumerate}
    \item $\displaystyle\lim_{n\to+\infty} \frac{3n^2-1}{n^2+4}$
    \item $\displaystyle\lim_{n\to+\infty} \frac{n+1}{\sqrt{n}}$
    \item $\displaystyle\lim_{n\to+\infty} \left(\frac{2}{3}\right)^n$
    \item $\displaystyle\lim_{n\to+\infty} n \cdot \left(\frac{1}{2}\right)^n$
  \end{enumerate}
  \end{multicols}
\end{exercice}

\begin{exercice}[5 : Suite r\'{e}currente]
  Soit $(u_n)$ d\'{e}finie par $u_0 = 1$ et $u_{n+1} = \sqrt{u_n + 2}$.
  \begin{enumerate}
    \item Montrer par r\'{e}currence que $u_n \in [0, 2]$ pour tout $n$.
    \item Montrer que $(u_n)$ est croissante.
    \item Conclure que $(u_n)$ converge et d\'{e}terminer sa limite.
  \end{enumerate}
\end{exercice}

\begin{exercice}[6 : Toile d'araign\'{e}e]
  Soit $(u_n)$ d\'{e}finie par $u_0 = 3$ et $u_{n+1} = \dfrac{1}{2}u_n + 1$.
  \begin{enumerate}
    \item R\'{e}soudre $f(x) = x$ pour trouver le point fixe.
    \item Calculer $u_1$, $u_2$, $u_3$.
    \item Repr\'{e}senter la toile d'araign\'{e}e.
    \item D\'{e}terminer la limite de $(u_n)$.
  \end{enumerate}
\end{exercice}
%==========================================================
%  CHAPITRE 5 — FONCTIONS EXPONENTIELLES
%==========================================================

\chapter{Fonctions Exponentielles}

\begin{rappel}[Logarithme n\'{e}p\'{e}rien : rappels essentiels]
  \begin{itemize}
    \item $\ln$ est d\'{e}finie sur $]0,+\infty[$ et est la r\'{e}ciproque de $\exp$.
    \item $\ln(e) = 1$, \quad $\ln(1) = 0$, \quad $\ln(ab) = \ln a + \ln b$.
    \item $(\ln x)' = \dfrac{1}{x}$ sur $]0,+\infty[$.
  \end{itemize}
\end{rappel}

%----------------------------------------------------------
\section{Fonction exponentielle n\'{e}p\'{e}rienne}
%----------------------------------------------------------

\subsection{D\'{e}finition et premi\`{e}res propri\'{e}t\'{e}s}

\begin{definition}[Fonction exponentielle n\'{e}p\'{e}rienne]
  La fonction \textbf{exponentielle n\'{e}p\'{e}rienne}, not\'{e}e $\exp$ ou $e^x$,
  est la \textbf{fonction r\'{e}ciproque} de $\ln$.

  \medskip
  On pose : $\exp(x) = e^x$, ce qui signifie :
  \[
    e^x = y \;\iff\; \ln(y) = x
    \qquad \bigl(x \in \R,\; y > 0\bigr)
  \]
\end{definition}

\begin{propriete}[Cons\'{e}quences directes]
  Pour tout $x \in \R$ et tout $y > 0$ :
  \[
    e^{\ln(y)} = y \qquad \ln(e^x) = x \qquad e^0 = 1 \qquad e^1 = e
  \]
\end{propriete}
\newpage

\begin{propriete}[R\`{e}gles de calcul]
  Pour tous r\'{e}els $x$, $y$ et tout rationnel $r$ :

  \bigskip
  \begin{center}
  \renewcommand{\arraystretch}{2}
  \begin{tabular}{|c|c|}
    \hline
    \textbf{\color{navy}Propri\'{e}t\'{e}}
      & \textbf{\color{navy}Formule} \\
    \hline\hline
    Produit
      & $e^x \times e^y = e^{x+y}$ \\
    \hline
    Quotient
      & $\dfrac{e^x}{e^y} = e^{x-y}$ \\
    \hline
    Inverse
      & $\dfrac{1}{e^x} = e^{-x}$ \\
    \hline
    Puissance
      & $\left(e^x\right)^r = e^{rx}$ \\
    \hline
    Comparaison
      & $e^x = e^y \iff x = y$ \quad et \quad $e^x > e^y \iff x > y$ \\
    \hline
    \'{E}quation
      & $e^x = y \iff x = \ln(y)$ \quad ($y > 0$) \\
    \hline
  \end{tabular}
  \end{center}
\end{propriete}

\begin{attention}
  \textbf{Erreurs fr\'{e}quentes \`{a} \'{e}viter :}
  \[
    e^{x+y} \neq e^x + e^y \qquad
    e^{x \times y} \neq e^x \times e^y \qquad
    (e^x)^y = e^{xy} \neq e^{e^{xy}}
  \]
\end{attention}

%----------------------------------------------------------
\subsection{Domaine de d\'{e}finition des fonctions avec \texorpdfstring{$e^{u(x)}$}{e\^{u(x)}}}
%----------------------------------------------------------

\begin{propriete}[Domaine de d\'{e}finition]
  \begin{center}
  \renewcommand{\arraystretch}{1.8}
  \begin{tabular}{|c|c|}
    \hline
    \textbf{\color{navy}Fonction}
      & \textbf{\color{navy}Domaine de d\'{e}finition} \\
    \hline\hline
    $f(x) = e^x$
      & $D_f = \R$ \\
    \hline
    $f(x) = e^{u(x)}$
      & $D_f = \{x \in \R \mid x \in D_u\}$ \\
    \hline
  \end{tabular}
  \end{center}

  \begin{remarque}
    $e^{u(x)}$ est toujours \textbf{strictement positif} :
    $e^{u(x)} > 0$ pour tout $x \in D_f$.
    Son domaine est simplement celui de $u$.
  \end{remarque}
\end{propriete}

%----------------------------------------------------------
\subsection{Limites usuelles}
%----------------------------------------------------------

\begin{propriete}[Limites de r\'{e}f\'{e}rence]
  \bigskip
  \begin{center}
  \renewcommand{\arraystretch}{2}
  \begin{tabular}{|c|c|}
    \hline
    \textbf{\color{navy}Limite}
      & \textbf{\color{navy}R\'{e}sultat} \\
    \hline\hline
    $\displaystyle\lim_{x \to +\infty} e^x$
      & $+\infty$ \\
    \hline
    $\displaystyle\lim_{x \to -\infty} e^x$
      & $0^+$ \\
    \hline
    $\displaystyle\lim_{x \to +\infty} \frac{e^x}{x^n}$
      & $+\infty$ \quad ($n \in \N^*$) \quad \textit{(croissance compar\'{e}e)} \\
    \hline
    $\displaystyle\lim_{x \to -\infty} x^n e^x$
      & $0$ \quad ($n \in \N^*$) \quad \textit{(croissance compar\'{e}e)} \\
    \hline
    $\displaystyle\lim_{x \to 0} \frac{e^x - 1}{x}$
      & $1$ \\
    \hline
  \end{tabular}
  \end{center}
\end{propriete}

\begin{remarque}
  \textbf{Croissance compar\'{e}e :} l'exponentielle cro\^{i}t
  \textbf{plus vite} que toute puissance de $x$.
  C'est fondamental pour les calculs de limites.
\end{remarque}

%----------------------------------------------------------
\subsection{Continuit\'{e} et d\'{e}rivabilit\'{e}}
%----------------------------------------------------------

\begin{propriete}[D\'{e}riv\'{e}e de $e^x$ et de $e^{u(x)}$]
  \begin{itemize}
    \item La fonction $x \mapsto e^x$ est d\'{e}rivable sur $\R$ et :
    \[
      \bigl(e^x\bigr)' = e^x
    \]
    \item Si $u$ est d\'{e}rivable sur $I$, alors $x \mapsto e^{u(x)}$
          est d\'{e}rivable sur $I$ et :
    \[
      \bigl(e^{u(x)}\bigr)' = u'(x) \cdot e^{u(x)}
    \]
  \end{itemize}
\end{propriete}

\begin{exemple}[R\'{e}solus : d\'{e}riv\'{e}es]
  \begin{itemize}
    \item $f(x) = e^{3x-1}$
    \quad $\Rightarrow$ \quad $u(x) = 3x-1$, $u'(x) = 3$
    \quad $\Rightarrow$ \quad $f'(x) = 3e^{3x-1}$

    \medskip
    \item $f(x) = e^{x^2}$
    \quad $\Rightarrow$ \quad $u(x) = x^2$, $u'(x) = 2x$
    \quad $\Rightarrow$ \quad $f'(x) = 2x\,e^{x^2}$

    \medskip
    \item $f(x) = x^2 e^x$
    \quad $\Rightarrow$ \quad r\`{e}gle du produit :
    $f'(x) = 2x\,e^x + x^2\,e^x = (2x+x^2)\,e^x = x(x+2)\,e^x$
  \end{itemize}
\end{exemple}

%----------------------------------------------------------
\subsection{Repr\'{e}sentation graphique}
%----------------------------------------------------------

\begin{propriete}[Courbe de $e^x$]
  La courbe de $f(x) = e^x$ passe par $(0,1)$, $(1,e)$ et $(-1, 1/e)$.
  Elle admet l'axe des abscisses comme \textbf{asymptote horizontale}
  en $-\infty$.
\end{propriete}

\begin{center}
\begin{tikzpicture}
  \begin{axis}[
    axis lines  = center,
    xlabel      = {$x$},
    ylabel      = {$y$},
    xmin = -3, xmax = 2.8,
    ymin = -0.5, ymax = 12,
    xtick       = {-3,-2,-1,0,1,2},
    ytick       = {0,2,4,6,8,10},
    grid        = major,
    grid style  = {gray!20},
    width = 10cm, height = 8cm,
  ]
    % Courbe e^x
    \addplot[navy, very thick, domain=-3:2.5, samples=150]
      {exp(x)};
    \addlegendentry{$f(x)=e^x$}

    % Courbe ln(x) (réciproque)
    \addplot[forest, thick, dashed, domain=0.05:2.8, samples=150]
      {ln(x)};
    \addlegendentry{$f^{-1}(x)=\ln x$}

    % Bissectrice y=x
    \addplot[burgundy, dashed, thin, domain=-2.5:2.8]
      {x};
    \addlegendentry{$y = x$}

    % Asymptote y=0
    \addplot[gray, dotted, domain=-3:0]
      {0};

    % Points remarquables
    \addplot[mark=*, mark size=3pt, navy]
      coordinates {(0,1)(1,2.718)};
    \node[above left, navy, font=\small] at (axis cs:0,1) {$(0,1)$};
    \node[right, navy, font=\small] at (axis cs:1,2.718) {$(1,e)$};

    \addplot[mark=*, mark size=3pt, forest]
      coordinates {(1,0)(2.718,1)};
    \node[below, forest, font=\small] at (axis cs:1,0) {$(1,0)$};
  \end{axis}
\end{tikzpicture}
\end{center}

%----------------------------------------------------------
\section{Fonction exponentielle de base \texorpdfstring{$a$}{a}}
%----------------------------------------------------------

\begin{definition}[Exponentielle de base $a$]
  Pour $a \in ]0,+\infty[$ et $a \neq 1$,
  la \textbf{fonction exponentielle de base $a$} est d\'{e}finie par :
  \[
    \forall x \in \R, \qquad a^x = e^{x\ln a}
  \]
\end{definition}

\begin{propriete}[R\`{e}gles de calcul]
  Pour tous $x, y \in \R$ et $r \in \Q$ :
  \[
    a^x \times a^y = a^{x+y} \qquad
    \frac{a^x}{a^y} = a^{x-y} \qquad
    \frac{1}{a^x} = a^{-x} \qquad
    (a^x)^r = a^{rx}
  \]
  \[
    a^x = e^{x\ln a} \qquad
    \log_a(a^x) = x \qquad
    a^{\log_a(x)} = x \quad (x > 0)
  \]
  \[
    a^x = a^y \iff x = y \qquad
    a^x = y \iff x = \log_a(y) \quad (y > 0)
  \]
\end{propriete}
\newpage
\begin{propriete}[Variations selon la valeur de $a$]
  \bigskip
  \begin{center}
  \renewcommand{\arraystretch}{2}
  \begin{tabular}{|p{6cm}|p{6cm}|}
    \hline
    \textbf{\color{burgundy}$0 < a < 1$ \quad (d\'{e}croissante)}
      & \textbf{\color{forest}$a > 1$ \quad (croissante)} \\
    \hline\hline
    $a^x > a^y \iff x < y$
      & $a^x > a^y \iff x > y$ \\
    \hline
    $\displaystyle\lim_{x \to +\infty} a^x = 0$
      & $\displaystyle\lim_{x \to +\infty} a^x = +\infty$ \\
    \hline
    $\displaystyle\lim_{x \to -\infty} a^x = +\infty$
      & $\displaystyle\lim_{x \to -\infty} a^x = 0$ \\
    \hline
    \multicolumn{2}{|c|}{%
      $\displaystyle\lim_{x \to 0} \dfrac{a^x - 1}{x} = \ln a$
    } \\
    \hline
    \multicolumn{2}{|c|}{%
      $(a^x)' = \ln(a) \cdot a^x$
    } \\
    \hline
  \end{tabular}
  \end{center}
\end{propriete}

\begin{center}
\begin{tikzpicture}
  \begin{axis}[
    axis lines  = center,
    xlabel      = {$x$},
    ylabel      = {$y$},
    xmin = -2.5, xmax = 2.5,
    ymin = -0.3, ymax = 7,
    xtick       = {-2,-1,0,1,2},
    ytick       = {0,1,2,3,4,5,6},
    grid        = major,
    grid style  = {gray!20},
    width = 10cm, height = 8cm,
  ]
    % a=2 (a > 1, croissante)
    \addplot[navy, very thick, domain=-2.4:2.4, samples=100]
      {2^x};
    \addlegendentry{$2^x$ \; ($a=2>1$, croissante)}

    % a=1/2 (0 < a < 1, décroissante)
    \addplot[burgundy, very thick, domain=-2.4:2.4, samples=100]
      {(0.5)^x};
    \addlegendentry{$(1/2)^x$ \; ($a=1/2<1$, d\'{e}croissante)}

    % a=e (courbe de référence)
    \addplot[forest, thick, dashed, domain=-2.4:1.9, samples=100]
      {exp(x)};
    \addlegendentry{$e^x$ \; (r\'{e}f\'{e}rence)}

    % Point commun (0,1)
    \addplot[mark=*, mark size=4pt, black]
      coordinates {(0,1)};
    \node[above right, black, font=\small\bfseries]
      at (axis cs:0,1) {$(0,1)$};
  \end{axis}
\end{tikzpicture}
\end{center}

\begin{exemple}[R\'{e}solus : \'{e}quations et in\'{e}quations avec $a^x$]
  \begin{itemize}
    \item R\'{e}soudre $2^x = 8$ :
    \quad $2^x = 2^3 \Rightarrow x = 3$

    \medskip
    \item R\'{e}soudre $3^x = 5$ :
    \quad $x\ln 3 = \ln 5 \Rightarrow x = \dfrac{\ln 5}{\ln 3}$

    \medskip
    \item R\'{e}soudre $\left(\dfrac{1}{2}\right)^x > 4$ :
    \quad $2^{-x} > 2^2 \Rightarrow -x > 2 \Rightarrow x < -2$
    \quad (in\'{e}galit\'{e} retourn\'{e}e car $a = 1/2 < 1$)
  \end{itemize}
\end{exemple}

%----------------------------------------------------------
\section{Signe et \'{e}quations avec \texorpdfstring{$e^{u(x)}$}{e\^{u(x)}}}
%----------------------------------------------------------

\begin{propriete}[Signe de $e^{u(x)}$]
  \[
    e^{u(x)} > 0 \quad \text{pour tout } x \in D_f
  \]
  On ne peut \textbf{jamais} avoir $e^{u(x)} = 0$ ou $e^{u(x)} < 0$.
\end{propriete}

\begin{methode}[R\'{e}soudre $e^{u(x)} = e^{v(x)}$ et $e^{u(x)} > e^{v(x)}$]
  \[
    e^{u(x)} = e^{v(x)} \iff u(x) = v(x)
  \]
  \[
    e^{u(x)} > e^{v(x)} \iff u(x) > v(x)
  \]
  car $x \mapsto e^x$ est \textbf{strictement croissante}.
\end{methode}

\begin{exemple}[R\'{e}solus : \'{e}quations avec $e^{u(x)}$]
  \begin{itemize}
    \item $e^{2x-1} = e^{x+3}$
    \quad $\Rightarrow$ \quad $2x-1 = x+3$
    \quad $\Rightarrow$ \quad $x = 4$

    \medskip
    \item $e^{x^2} = e^{2x}$
    \quad $\Rightarrow$ \quad $x^2 = 2x$
    \quad $\Rightarrow$ \quad $x(x-2) = 0$
    \quad $\Rightarrow$ \quad $x = 0$ ou $x = 2$

    \medskip
    \item $e^{x^2 - 3} > 1 = e^0$
    \quad $\Rightarrow$ \quad $x^2 - 3 > 0$
    \quad $\Rightarrow$ \quad $x < -\sqrt{3}$ ou $x > \sqrt{3}$
  \end{itemize}
\end{exemple}

\newpage
%----------------------------------------------------------
\section*{Exercices du chapitre}
\addcontentsline{toc}{section}{Exercices : Fonctions exponentielles}
%----------------------------------------------------------

\begin{exercice}[1 : Calculs avec $e^x$]
  Simplifier les expressions suivantes :
  \begin{multicols}{2}
  \begin{enumerate}
    \item $e^3 \times e^{-1}$
    \item $\dfrac{e^5}{e^2}$
    \item $(e^2)^3$
    \item $e^{\ln 5}$
    \item $\ln(e^{-4})$
    \item $e^{2\ln 3}$
  \end{enumerate}
  \end{multicols}
\end{exercice}

\begin{exercice}[2 : D\'{e}riv\'{e}es]
  Calculer la d\'{e}riv\'{e}e des fonctions suivantes :
  \begin{multicols}{2}
  \begin{enumerate}
    \item $f(x) = e^{4x+1}$
    \item $f(x) = e^{x^2-3}$
    \item $f(x) = (x+1)e^x$
    \item $f(x) = \dfrac{e^x}{x}$
    \item $f(x) = e^{\sqrt{x}}$
    \item $f(x) = e^{-x^2}$
  \end{enumerate}
  \end{multicols}
\end{exercice}

\begin{exercice}[3 : \'{E}quations et in\'{e}quations]
  R\'{e}soudre dans $\R$ :
  \begin{enumerate}
    \item $e^{3x} = e^{x+4}$
    \item $e^{x^2-1} = 1$
    \item $e^{2x} - 3e^x + 2 = 0$
          \quad (poser $X = e^x$)
    \item $e^x > e^{2-x}$
    \item $2^x = 7$ \quad (donner le r\'{e}sultat avec $\ln$)
  \end{enumerate}
\end{exercice}

\begin{exercice}[4 : Limites]
  Calculer les limites suivantes :
  \begin{multicols}{2}
  \begin{enumerate}
    \item $\displaystyle\lim_{x\to+\infty} \frac{e^x}{x^3}$
    \item $\displaystyle\lim_{x\to-\infty} x^2 e^x$
    \item $\displaystyle\lim_{x\to+\infty} (e^x - x^{100})$
    \item $\displaystyle\lim_{x\to 0} \frac{e^{2x}-1}{x}$
  \end{enumerate}
  \end{multicols}
\end{exercice}

\begin{exercice}[5 : \'{E}tude compl\`{e}te]
  Soit $f(x) = (x-1)e^x$ d\'{e}finie sur $\R$.
  \begin{enumerate}
    \item Calculer les limites de $f$ en $\pm\infty$.
    \item Calculer $f'(x)$ et dresser le tableau de variations.
    \item D\'{e}terminer les points d'inflexion (calculer $f''$).
    \item Tracer la courbe $\mathcal{C}_f$.
  \end{enumerate}
\end{exercice}
%==========================================================
%  CHAPITRE 7 — FONCTIONS PRIMITIVES
%==========================================================

\chapter{Fonctions Primitives}

\begin{rappel}[Lien avec la d\'{e}rivation]
  Calculer une primitive est l'\textbf{op\'{e}ration inverse} de la d\'{e}rivation.
  Si $F$ est une primitive de $f$, alors $F' = f$.
  \begin{itemize}
    \item D\'{e}river : $F \xrightarrow{\ '\ } f$
    \item Primitiver : $f \xrightarrow{\int} F$
  \end{itemize}
\end{rappel}

%----------------------------------------------------------
\section{D\'{e}finition et propri\'{e}t\'{e}s g\'{e}n\'{e}rales}
%----------------------------------------------------------

\begin{definition}[Fonction primitive]
  Soit $f$ une fonction d\'{e}finie sur un intervalle $I$.
  On dit que $F$ est une \textbf{primitive de $f$ sur $I$} si :
  \begin{itemize}
    \item $F$ est d\'{e}rivable sur $I$
    \item $\forall x \in I, \quad F'(x) = f(x)$
  \end{itemize}
\end{definition}

\begin{theoreme}[Existence et unicit\'{e}]
  \begin{itemize}
    \item Toute fonction \textbf{continue} sur un intervalle $I$ admet
          des primitives sur $I$.
    \item Si $F$ est une primitive de $f$ sur $I$, alors toute fonction
          $G$ d\'{e}finie par $G(x) = F(x) + k$ \; ($k \in \R$) est aussi
          une primitive de $f$ sur $I$.
    \item \textbf{Unicit\'{e} avec condition initiale :} si $x_0 \in I$ et
          $y_0 \in \R$, il existe \textbf{une seule} primitive $F$ de $f$ sur $I$
          v\'{e}rifiant $F(x_0) = y_0$.
  \end{itemize}
\end{theoreme}

\begin{remarque}
  Deux primitives d'une m\^{e}me fonction diff\`{e}rent d'une constante.
  On \'{e}crit souvent $F(x) + k$ ou $F(x) + C$ pour d\'{e}signer
  \textbf{la famille des primitives} de $f$.
\end{remarque}

\begin{propriete}[Lin\'{e}arit\'{e} des primitives]
  Si $F$ est une primitive de $f$ et $G$ une primitive de $g$ sur $I$,
  et si $k \in \R$, alors :
  \begin{itemize}
    \item $F + G$ est une primitive de $f + g$ sur $I$
    \item $kF$ est une primitive de $kf$ sur $I$
  \end{itemize}
  \[
    \text{En r\'{e}sum\'{e} :} \quad
    \int \bigl[\alpha f(x) + \beta g(x)\bigr]\,dx
    = \alpha F(x) + \beta G(x) + C
  \]
\end{propriete}

\begin{methode}[Trouver la primitive v\'{e}rifiant une condition initiale]
  \begin{enumerate}
    \item \'{E}crire la forme g\'{e}n\'{e}rale $F(x) = \ldots + k$.
    \item Utiliser la condition $F(x_0) = y_0$ pour trouver $k$.
    \item Substituer la valeur de $k$.
  \end{enumerate}
\end{methode}

\begin{exemple}[R\'{e}solu : condition initiale]
  Trouver la primitive de $f(x) = 2x + 3$ qui v\'{e}rifie $F(1) = 5$.

  \medskip
  Forme g\'{e}n\'{e}rale : $F(x) = x^2 + 3x + k$.

  $F(1) = 5 \Rightarrow 1 + 3 + k = 5 \Rightarrow k = 1$.

  Donc $F(x) = x^2 + 3x + 1$.
\end{exemple}

%----------------------------------------------------------
\section{Formulaire des primitives usuelles}
%----------------------------------------------------------

\begin{propriete}[Primitives des fonctions de base]
  \bigskip
  \begin{center}
  \renewcommand{\arraystretch}{2.4}
  \begin{tabular}{|c|c|c|}
    \hline
    \textbf{\color{navy}Fonction $f(x)$}
      & \textbf{\color{navy}Primitive $F(x)$}
      & \textbf{\color{navy}Domaine} \\
    \hline\hline
    $a$ \; (constante)
      & $ax$
      & $\R$ \\
    \hline
    $x$
      & $\dfrac{x^2}{2}$
      & $\R$ \\
    \hline
    $x^n$ \; ($n \in \N$)
      & $\dfrac{x^{n+1}}{n+1}$
      & $\R$ \\
    \hline
    $x^n$ \; ($n \in \Z,\; n \neq -1$)
      & $\dfrac{x^{n+1}}{n+1}$
      & $\R^*$ ou $]0,+\infty[$ \\
    \hline
    $\dfrac{1}{x^2}$
      & $-\dfrac{1}{x}$
      & $\R^*$ \\
    \hline
    $\dfrac{1}{\sqrt{x}}$
      & $2\sqrt{x}$
      & $]0,+\infty[$ \\
    \hline
    $\dfrac{1}{x}$
      & $\ln|x|$
      & $\R^*$ \\
    \hline
    $e^x$
      & $e^x$
      & $\R$ \\
    \hline
    $\cos x$
      & $\sin x$
      & $\R$ \\
    \hline
    $\sin x$
      & $-\cos x$
      & $\R$ \\
    \hline
  \end{tabular}
  \end{center}
\end{propriete}

%----------------------------------------------------------
\section{Primitives des fonctions compos\'{e}es}
%----------------------------------------------------------

\begin{propriete}[Formulaire des primitives compos\'{e}es]
  Soit $u$ une fonction d\'{e}rivable sur $I$.

  \bigskip
  \begin{center}
  \renewcommand{\arraystretch}{2.4}
  \begin{tabular}{|c|c|c|}
    \hline
    \textbf{\color{forest}Fonction $f(x)$}
      & \textbf{\color{forest}Primitive $F(x)$}
      & \textbf{\color{forest}Condition} \\
    \hline\hline
    $u'(x) \cdot u(x)$
      & $\dfrac{[u(x)]^2}{2}$
      & \\
    \hline
    $u'(x) \cdot [u(x)]^n$ \; ($n \neq -1$)
      & $\dfrac{[u(x)]^{n+1}}{n+1}$
      & $n \in \Z^*\setminus\{-1\}$ \\
    \hline
    $\dfrac{u'(x)}{u(x)}$
      & $\ln|u(x)|$
      & $u(x) \neq 0$ \\
    \hline
    $\dfrac{u'(x)}{\sqrt{u(x)}}$
      & $2\sqrt{u(x)}$
      & $u(x) > 0$ \\
    \hline
    $u'(x) \cdot e^{u(x)}$
      & $e^{u(x)}$
      & \\
    \hline
    $u'(x) \cdot \cos(u(x))$
      & $\sin(u(x))$
      & \\
    \hline
    $u'(x) \cdot \sin(u(x))$
      & $-\cos(u(x))$
      & \\
    \hline
  \end{tabular}
  \end{center}
\end{propriete}

\begin{attention}
  Pour reconna\^{i}tre une primitive compos\'{e}e, il faut \textbf{identifier $u$ et $u'$}
  dans l'expression. Si $u'$ n'appara\^{i}t pas exactement, on peut
  \textbf{ajuster par une constante} multiplicative.
\end{attention}

\begin{methode}[Calculer une primitive compos\'{e}e : ajustement par constante]
  \begin{enumerate}
    \item \textbf{Identifier} $u(x)$ et calculer $u'(x)$.
    \item \textbf{V\'{e}rifier} si l'expression contient $u'(x)$ (exactement
          ou \`{a} une constante pr\`{e}s).
    \item \textbf{Ajuster} en multipliant/divisant par la constante
          n\'{e}cessaire.
    \item \textbf{V\'{e}rifier} en d\'{e}rivant le r\'{e}sultat.
  \end{enumerate}
\end{methode}

\begin{exemple}[R\'{e}solus : primitives compos\'{e}es]
  \begin{itemize}

    \item $f(x) = (2x+1)^4$

    $u = 2x+1$, $u' = 2$.
    On voit $u^4$ mais il manque $u'=2$ devant.
    On \'{e}crit : $f(x) = \dfrac{1}{2} \cdot 2 \cdot (2x+1)^4 = \dfrac{1}{2} \cdot u' \cdot u^4$.

    Donc $F(x) = \dfrac{1}{2} \cdot \dfrac{(2x+1)^5}{5} = \dfrac{(2x+1)^5}{10}$.

    \textit{V\'{e}rif. :} $F'(x) = \dfrac{5 \cdot 2(2x+1)^4}{10} = (2x+1)^4$ $\checkmark$

    \bigskip
    \item $f(x) = \dfrac{3x^2}{x^3 + 1}$

    $u = x^3+1$, $u' = 3x^2$.
    On reconna\^{i}t $\dfrac{u'}{u}$.
    Donc $F(x) = \ln|x^3+1|$.

    \bigskip
    \item $f(x) = (4x-3)\,e^{2x^2 - 3x}$

    $u = 2x^2-3x$, $u' = 4x-3$.
    On reconna\^{i}t $u' \cdot e^u$.
    Donc $F(x) = e^{2x^2 - 3x}$.

    \bigskip
    \item $f(x) = \dfrac{x}{\sqrt{x^2+1}}$

    $u = x^2+1$, $u' = 2x$.
    On voit $\dfrac{u'}{2\sqrt{u}}$, donc :
    $f(x) = \dfrac{1}{2} \cdot \dfrac{2x}{\sqrt{x^2+1}} = \dfrac{1}{2} \cdot \dfrac{u'}{\sqrt{u}}$.
    Donc $F(x) = \dfrac{1}{2} \cdot 2\sqrt{x^2+1} = \sqrt{x^2+1}$.

  \end{itemize}
\end{exemple}

%----------------------------------------------------------
\section{Primitives et tableaux de variations}
%----------------------------------------------------------

\begin{propriete}[Signe de $f$ et monotonie de $F$]
  Puisque $F' = f$ :
  \begin{itemize}
    \item $f(x) > 0$ sur $I$ $\Rightarrow$ $F$ est \textbf{croissante} sur $I$
    \item $f(x) < 0$ sur $I$ $\Rightarrow$ $F$ est \textbf{d\'{e}croissante} sur $I$
    \item $f(x_0) = 0$ et $f$ change de signe en $x_0$
          $\Rightarrow$ $F$ admet un \textbf{extremum} en $x_0$
  \end{itemize}
\end{propriete}

\newpage
%----------------------------------------------------------
\section*{Exercices du chapitre}
\addcontentsline{toc}{section}{Exercices : Fonctions primitives}
%----------------------------------------------------------

\begin{exercice}[1 : Primitives imm\'{e}diates]
  Donner une primitive de chacune des fonctions suivantes :
  \begin{multicols}{2}
  \begin{enumerate}
    \item $f(x) = 5x^3 - 2x + 7$
    \item $f(x) = \dfrac{3}{x^2}$
    \item $f(x) = \sqrt{x} + \dfrac{1}{\sqrt{x}}$
    \item $f(x) = 4e^x - \dfrac{2}{x}$
    \item $f(x) = \cos x - \sin x$
    \item $f(x) = \dfrac{1}{x} + x^4$
  \end{enumerate}
  \end{multicols}
\end{exercice}

\begin{exercice}[2 : Primitives compos\'{e}es]
  Reconnaître et calculer une primitive de :
  \begin{multicols}{2}
  \begin{enumerate}
    \item $f(x) = (3x-1)^5$
    \item $f(x) = \dfrac{2x}{x^2+4}$
    \item $f(x) = x\,e^{x^2}$
    \item $f(x) = \dfrac{1}{(2x+1)^3}$
    \item $f(x) = \dfrac{x}{\sqrt{1-x^2}}$
    \item $f(x) = (2x-1)e^{x^2-x}$
  \end{enumerate}
  \end{multicols}
\end{exercice}

\begin{exercice}[3 : Condition initiale]
  Trouver la primitive de $f$ v\'{e}rifiant la condition donn\'{e}e :
  \begin{enumerate}
    \item $f(x) = 3x^2 - 4x$ et $F(0) = 2$
    \item $f(x) = e^x + 1$ et $F(0) = 0$
    \item $f(x) = \dfrac{1}{x}$ et $F(1) = 3$
    \item $f(x) = \cos x$ et $F\!\left(\dfrac{\pi}{2}\right) = 1$
  \end{enumerate}
\end{exercice}

\begin{exercice}[4 : Lin\'{e}arit\'{e}]
  Calculer une primitive de :
  \begin{enumerate}
    \item $f(x) = 3(2x+1)^4 - \dfrac{5}{x}$ \quad sur $]0,+\infty[$
    \item $f(x) = 2e^{3x} + 4\cos x$
    \item $f(x) = \dfrac{x^3-x+1}{x^2}$ \quad (simplifier d'abord)
  \end{enumerate}
\end{exercice}

\begin{exercice}[5 : V\'{e}rification par d\'{e}rivation]
  V\'{e}rifier que $F$ est bien une primitive de $f$ :
  \begin{enumerate}
    \item $f(x) = \dfrac{x}{\sqrt{x^2+1}}$ \quad et \quad $F(x) = \sqrt{x^2+1}$
    \item $f(x) = \ln x$ \quad et \quad $F(x) = x\ln x - x$
    \item $f(x) = \sin^2 x$ \quad et \quad $F(x) = \dfrac{x}{2} - \dfrac{\sin(2x)}{4}$
  \end{enumerate}
\end{exercice}
%==========================================================
%  CHAPITRE 6 — FONCTIONS LOGARITHMIQUES
%==========================================================

\chapter{Fonctions Logarithmiques}

\begin{rappel}[Lien avec la fonction exponentielle]
  \begin{itemize}
    \item $\ln$ et $\exp$ sont \textbf{r\'{e}ciproques} l'une de l'autre :
          $e^{\ln x} = x$ et $\ln(e^x) = x$.
    \item $e^x > 0$ pour tout $x \in \R$.
    \item $\ln$ est d\'{e}finie uniquement sur $]0,+\infty[$.
  \end{itemize}
\end{rappel}

%----------------------------------------------------------
\section{Fonction Logarithme N\'{e}p\'{e}rien}
%----------------------------------------------------------

\subsection{D\'{e}finition}

\begin{definition}[Logarithme n\'{e}p\'{e}rien]
  La \textbf{fonction logarithme n\'{e}p\'{e}rien}, not\'{e}e $\ln$,
  est l'unique primitive de $x \mapsto \dfrac{1}{x}$ sur $]0,+\infty[$
  qui s'annule en $1$ :
  \[
    (\ln x)' = \frac{1}{x} \quad \text{sur } ]0,+\infty[
    \qquad \text{et} \qquad \ln(1) = 0
  \]
\end{definition}

\begin{propriete}[Cons\'{e}quences directes]
  \[
    \ln(1) = 0 \qquad \ln(e) = 1 \qquad
    \ln(x) = 0 \iff x = 1 \qquad
    \ln(x) > 0 \iff x > 1
  \]
\end{propriete}

%----------------------------------------------------------
\subsection{R\`{e}gles de calcul}
%----------------------------------------------------------

\begin{propriete}[Propri\'{e}t\'{e}s alg\'{e}briques]
  Pour tous $x > 0$, $y > 0$ et $r \in \Q$ :

  \bigskip
  \begin{center}
  \renewcommand{\arraystretch}{2}
  \begin{tabular}{|c|c|}
    \hline
    \textbf{\color{navy}Propri\'{e}t\'{e}}
      & \textbf{\color{navy}Formule} \\
    \hline\hline
    Produit
      & $\ln(xy) = \ln x + \ln y$ \\
    \hline
    Quotient
      & $\ln\!\left(\dfrac{x}{y}\right) = \ln x - \ln y$ \\
    \hline
    Inverse
      & $\ln\!\left(\dfrac{1}{x}\right) = -\ln x$ \\
    \hline
    Puissance
      & $\ln(x^r) = r\ln x$ \\
    \hline
    Racine carr\'{e}e
      & $\ln(\sqrt{x}) = \dfrac{1}{2}\ln x$ \\
    \hline
  \end{tabular}
  \end{center}
\end{propriete}

\begin{propriete}[Comparaison et \'{e}quations]
  Pour $x > 0$ et $y > 0$ :
  \[
    \ln x = \ln y \iff x = y \qquad
    \ln x < \ln y \iff x < y
  \]
  \[
    \ln x = y \iff x = e^y \qquad
    \ln x = r \iff x = e^r \quad (r \in \R)
  \]
\end{propriete}

\begin{attention}
  \textbf{Erreurs fr\'{e}quentes \`{a} \'{e}viter :}
  \[
    \ln(x + y) \neq \ln x + \ln y \qquad
    \ln(x \times y) \neq \ln x \times \ln y \qquad
    \frac{\ln x}{\ln y} \neq \ln\!\left(\frac{x}{y}\right)
  \]
\end{attention}

%----------------------------------------------------------
\subsection{Domaine de d\'{e}finition}
%----------------------------------------------------------

\begin{propriete}[Domaine des fonctions avec $\ln$]
  \bigskip
  \begin{center}
  \renewcommand{\arraystretch}{2}
  \begin{tabular}{|c|c|}
    \hline
    \textbf{\color{navy}Fonction}
      & \textbf{\color{navy}Domaine de d\'{e}finition} \\
    \hline\hline
    $f(x) = \ln x$
      & $D_f = \;]0,+\infty[$ \\
    \hline
    $f(x) = \ln\bigl(u(x)\bigr)$
      & $D_f = \{x \in D_u \mid u(x) > 0\}$ \\
    \hline
    $f(x) = \ln\bigl(u(x)^2\bigr)$
      & $D_f = \{x \in D_u \mid u(x) \neq 0\}$ \\
    \hline
    $f(x) = \ln\bigl|u(x)\bigr|$
      & $D_f = \{x \in D_u \mid u(x) \neq 0\}$ \\
    \hline
  \end{tabular}
  \end{center}
\end{propriete}

\begin{exemple}[R\'{e}solus : domaines de d\'{e}finition]
  \begin{itemize}
    \item $f(x) = \ln(2x - 3)$ :
    \quad $2x - 3 > 0 \Rightarrow x > \dfrac{3}{2}$
    \quad $\Rightarrow D_f = \left]\dfrac{3}{2},+\infty\right[$

    \medskip
    \item $f(x) = \ln(x^2 - 4)$ :
    \quad $x^2 - 4 > 0 \Rightarrow (x-2)(x+2) > 0$
    \quad $\Rightarrow D_f = \;]-\infty,-2[\;\cup\;]2,+\infty[$

    \medskip
    \item $f(x) = \ln(x^2 + 1)$ :
    \quad $x^2 + 1 > 0$ toujours vrai
    \quad $\Rightarrow D_f = \R$
  \end{itemize}
\end{exemple}

%----------------------------------------------------------
\subsection{Limites usuelles}
%----------------------------------------------------------

\begin{propriete}[Limites de r\'{e}f\'{e}rence]
  \bigskip
  \begin{center}
  \renewcommand{\arraystretch}{2}
  \begin{tabular}{|c|c|}
    \hline
    \textbf{\color{navy}Limite}
      & \textbf{\color{navy}R\'{e}sultat} \\
    \hline\hline
    $\displaystyle\lim_{x \to +\infty} \ln x$
      & $+\infty$ \\
    \hline
    $\displaystyle\lim_{x \to 0^+} \ln x$
      & $-\infty$ \\
    \hline
    $\displaystyle\lim_{x \to +\infty} \dfrac{\ln x}{x^n}$
      & $0$ \quad ($n \in \N^*$) \quad \textit{(croissance compar\'{e}e)} \\
    \hline
    $\displaystyle\lim_{x \to 0^+} x^n \ln x$
      & $0$ \quad ($n \in \N^*$) \quad \textit{(croissance compar\'{e}e)} \\
    \hline
    $\displaystyle\lim_{x \to 1} \dfrac{\ln x}{x - 1}$
      & $1$ \\
    \hline
    $\displaystyle\lim_{x \to 0} \dfrac{\ln(1+x)}{x}$
      & $1$ \\
    \hline
  \end{tabular}
  \end{center}
\end{propriete}

%----------------------------------------------------------
\subsection{D\'{e}rivabilit\'{e}}
%----------------------------------------------------------

\begin{propriete}[D\'{e}riv\'{e}e de $\ln x$ et de $\ln u(x)$]
  \begin{itemize}
    \item La fonction $\ln$ est d\'{e}rivable sur $]0,+\infty[$ et :
    \[
      (\ln x)' = \frac{1}{x}
    \]
    \item Si $u$ est d\'{e}rivable et $u > 0$ sur $I$, alors :
    \[
      \bigl(\ln u(x)\bigr)' = \frac{u'(x)}{u(x)}
    \]
  \end{itemize}
\end{propriete}

\begin{exemple}[R\'{e}solus : d\'{e}riv\'{e}es]
  \begin{itemize}
    \item $f(x) = \ln(3x+1)$ :
    \quad $u = 3x+1$, $u' = 3$
    \quad $\Rightarrow f'(x) = \dfrac{3}{3x+1}$

    \medskip
    \item $f(x) = \ln(x^2+1)$ :
    \quad $u = x^2+1$, $u' = 2x$
    \quad $\Rightarrow f'(x) = \dfrac{2x}{x^2+1}$

    \medskip
    \item $f(x) = x\ln x$ :
    \quad r\`{e}gle du produit :
    $f'(x) = 1 \cdot \ln x + x \cdot \dfrac{1}{x} = \ln x + 1$

    \medskip
    \item $f(x) = \ln\!\left(\dfrac{x+1}{x-1}\right) = \ln(x+1) - \ln(x-1)$ :
    \quad $f'(x) = \dfrac{1}{x+1} - \dfrac{1}{x-1} = \dfrac{-2}{x^2-1}$
  \end{itemize}
\end{exemple}

%----------------------------------------------------------
\subsection{Signe et repr\'{e}sentation graphique}
%----------------------------------------------------------

\begin{propriete}[Signe de $\ln x$]
  \bigskip
  \begin{center}
  \renewcommand{\arraystretch}{1.8}
  \begin{tabular}{|c|ccccc|}
    \hline
    $x$      & $0$ && $1$ && $+\infty$ \\
    \hline
    $\ln x$  & \multicolumn{2}{c}{$-\infty \quad -$} & $0$ & $+$ & $+\infty$ \\
    \hline
  \end{tabular}
  \end{center}

  \medskip
  $\ln x < 0$ pour $0 < x < 1$ \quad et \quad $\ln x > 0$ pour $x > 1$.
\end{propriete}

\begin{center}
\begin{tikzpicture}
  \begin{axis}[
    axis lines  = center,
    xlabel      = {$x$},
    ylabel      = {$y$},
    xmin = -0.5, xmax = 6,
    ymin = -3,   ymax = 3,
    xtick       = {0,1,2,3,4,5},
    ytick       = {-3,-2,-1,0,1,2,3},
    grid        = major,
    grid style  = {gray!20},
    width = 10cm, height = 8cm,
  ]
    % Courbe ln(x)
    \addplot[navy, very thick, domain=0.05:6, samples=200]
      {ln(x)};
    \addlegendentry{$f(x)=\ln x$}

    % Courbe e^x (réciproque)
    \addplot[forest, thick, dashed, domain=-2.5:1.8, samples=150]
      {exp(x)};
    \addlegendentry{$f^{-1}(x)=e^x$}

    % Bissectrice y=x
    \addplot[burgundy, thin, dashed, domain=-0.5:3]
      {x};
    \addlegendentry{$y=x$}

    % Asymptote x=0
    \addplot[gray, dotted, domain=-3:3]
      coordinates {(0,-3)(0,3)};

    % Points remarquables
    \addplot[mark=*, mark size=3pt, navy]
      coordinates {(1,0)(2.718,1)};
    \node[below right, navy, font=\small] at (axis cs:1,0) {$(1,0)$};
    \node[right, navy, font=\small] at (axis cs:2.718,1) {$(e,1)$};

    % Zone négative
    \node[burgundy, font=\small\bfseries] at (axis cs:0.4,-1.5) {$\ln x < 0$};
    % Zone positive
    \node[forest, font=\small\bfseries] at (axis cs:4,0.8) {$\ln x > 0$};
  \end{axis}
\end{tikzpicture}
\end{center}

%----------------------------------------------------------
\section{Fonction Logarithme de base $a$}
%----------------------------------------------------------

\begin{definition}[Logarithme de base $a$]
  Pour $a \in \R^*_+$ et $a \neq 1$,
  la \textbf{fonction logarithme de base $a$} est d\'{e}finie par :
  \[
    \forall x \in\, ]0,+\infty[, \qquad
    \log_a(x) = \frac{\ln x}{\ln a}
  \]
  \textbf{Cas particulier :} si $a = 10$, $\log_{10}$ est le
  \textbf{logarithme d\'{e}cimal}, not\'{e} $\log$.
\end{definition}

\begin{propriete}[R\`{e}gles de calcul]
  Pour $x > 0$, $y > 0$, $r \in \Q$ :
  \[
    \log_a(xy) = \log_a x + \log_a y \qquad
    \log_a\!\left(\frac{x}{y}\right) = \log_a x - \log_a y
  \]
  \[
    \log_a(x^r) = r\log_a x \qquad
    \log_a 1 = 0 \qquad
    \log_a a = 1
  \]
  \[
    \log_a x = \log_a y \iff x = y \qquad
    \log_a x = r \iff x = a^r
  \]
\end{propriete}

\begin{propriete}[Variations et limites selon $a$]
  \bigskip
  \begin{center}
  \renewcommand{\arraystretch}{2}
  \begin{tabular}{|p{6cm}|p{6cm}|}
    \hline
    \textbf{\color{burgundy}$0 < a < 1$ \quad (d\'{e}croissante)}
      & \textbf{\color{forest}$a > 1$ \quad (croissante)} \\
    \hline\hline
    $\log_a x > \log_a y \iff x < y$
      & $\log_a x > \log_a y \iff x > y$ \\
    \hline
    $\displaystyle\lim_{x\to+\infty} \log_a x = -\infty$
      & $\displaystyle\lim_{x\to+\infty} \log_a x = +\infty$ \\
    \hline
    $\displaystyle\lim_{x\to 0^+} \log_a x = +\infty$
      & $\displaystyle\lim_{x\to 0^+} \log_a x = -\infty$ \\
    \hline
    \multicolumn{2}{|c|}{%
      $\displaystyle\lim_{x\to 0} \dfrac{a^x - 1}{x} = \ln a$
      \qquad
      $\bigl(\log_a x\bigr)' = \dfrac{1}{x \ln a}$
    } \\
    \hline
  \end{tabular}
  \end{center}
\end{propriete}

\begin{center}
\begin{tikzpicture}
  \begin{axis}[
    axis lines  = center,
    xlabel      = {$x$},
    ylabel      = {$y$},
    xmin = -0.3, xmax = 6,
    ymin = -3,   ymax = 3,
    xtick       = {0,1,2,3,4,5},
    ytick       = {-3,-2,-1,0,1,2},
    grid        = major,
    grid style  = {gray!20},
    width = 10cm, height = 8cm,
  ]
    % log base 2 (a>1, croissante)
    \addplot[navy, very thick, domain=0.05:6, samples=200]
      {ln(x)/ln(2)};
    \addlegendentry{$\log_2 x$ \; ($a=2>1$)}

    % log base 1/2 (0<a<1, décroissante)
    \addplot[burgundy, very thick, domain=0.05:6, samples=200]
      {ln(x)/ln(0.5)};
    \addlegendentry{$\log_{1/2} x$ \; ($a=1/2<1$)}

    % ln (référence)
    \addplot[forest, thick, dashed, domain=0.05:6, samples=200]
      {ln(x)};
    \addlegendentry{$\ln x$ \; (r\'{e}f\'{e}rence, $a=e$)}

    % Point commun (1,0)
    \addplot[mark=*, mark size=4pt, black]
      coordinates {(1,0)};
    \node[above right, black, font=\small\bfseries]
      at (axis cs:1,0) {$(1,0)$};
  \end{axis}
\end{tikzpicture}
\end{center}

\begin{exemple}[R\'{e}solus : logarithme de base $a$]
  \begin{itemize}
    \item $\log_2(32) = \log_2(2^5) = 5$
    \medskip
    \item $\log_3(x) = 2 \Rightarrow x = 3^2 = 9$
    \medskip
    \item $\log_{10}(1000) = \log_{10}(10^3) = 3$
    \medskip
    \item R\'{e}soudre $\log_2 x > 3$ :
    \quad $x > 2^3 = 8$ \quad (car $a = 2 > 1$, $\log_2$ croissante)
  \end{itemize}
\end{exemple}

\newpage
%----------------------------------------------------------
\section*{Exercices du chapitre}
\addcontentsline{toc}{section}{Exercices : Fonctions logarithmiques}
%----------------------------------------------------------

\begin{exercice}[1 : Calculs avec $\ln$]
  Simplifier les expressions suivantes :
  \begin{multicols}{2}
  \begin{enumerate}
    \item $\ln(e^5)$
    \item $\ln\!\left(\dfrac{1}{e^3}\right)$
    \item $e^{2\ln 3}$
    \item $\ln(\sqrt{e})$
    \item $\ln(2) + \ln(3) - \ln(6)$
    \item $3\ln(2) - \ln(4)$
  \end{enumerate}
  \end{multicols}
\end{exercice}

\begin{exercice}[2 : Domaine de d\'{e}finition]
  D\'{e}terminer le domaine de d\'{e}finition de :
  \begin{enumerate}
    \item $f(x) = \ln(5x - 2)$
    \item $f(x) = \ln(x^2 - 9)$
    \item $f(x) = \ln(x^2 + 2x + 1)$
    \item $f(x) = \ln(x) + \ln(1-x)$
  \end{enumerate}
\end{exercice}

\begin{exercice}[3 : D\'{e}riv\'{e}es]
  Calculer la d\'{e}riv\'{e}e des fonctions suivantes :
  \begin{multicols}{2}
  \begin{enumerate}
    \item $f(x) = \ln(x^2 + 3)$
    \item $f(x) = x^2 \ln x$
    \item $f(x) = \dfrac{\ln x}{x}$
    \item $f(x) = \ln\!\left(\dfrac{x+2}{x-1}\right)$
    \item $f(x) = \ln(\cos x)$
    \item $f(x) = \ln(x + \sqrt{x^2+1})$
  \end{enumerate}
  \end{multicols}
\end{exercice}

\begin{exercice}[4 : \'{E}quations et in\'{e}quations]
  R\'{e}soudre dans $\R$ :
  \begin{enumerate}
    \item $\ln(2x+1) = \ln(x+3)$
    \item $\ln(x^2) = \ln(4)$
    \item $\ln(x) + \ln(x-2) = \ln(3)$
    \item $\ln^2(x) - 3\ln(x) + 2 = 0$
          \quad (poser $X = \ln x$)
    \item $\ln(x) > \ln(3x - 2)$
  \end{enumerate}
\end{exercice}

\begin{exercice}[5 : Limites]
  Calculer les limites suivantes :
  \begin{multicols}{2}
  \begin{enumerate}
    \item $\displaystyle\lim_{x\to+\infty} \frac{\ln x}{x}$
    \item $\displaystyle\lim_{x\to 0^+} x \ln x$
    \item $\displaystyle\lim_{x\to+\infty} \frac{\ln x}{\sqrt{x}}$
    \item $\displaystyle\lim_{x\to 0} \frac{\ln(1+3x)}{x}$
  \end{enumerate}
  \end{multicols}
\end{exercice}

\begin{exercice}[6 : \'{E}tude compl\`{e}te]
  Soit $f(x) = x - \ln x$ sur $]0,+\infty[$.
  \begin{enumerate}
    \item Calculer les limites de $f$ en $0^+$ et en $+\infty$.
    \item Calculer $f'(x)$ et dresser le tableau de variations.
    \item D\'{e}terminer le minimum de $f$.
    \item Montrer que $\ln x \leq x - 1$ pour tout $x > 0$.
  \end{enumerate}
\end{exercice}
%==========================================================
%  CHAPITRE 8 — CALCUL INTÉGRAL
%==========================================================

\chapter{Calcul Int\'{e}gral}

\begin{rappel}[Lien avec les primitives]
  Si $F$ est une primitive de $f$ sur $[a,b]$, alors :
  \[
    \int_a^b f(x)\,dx = \bigl[F(x)\bigr]_a^b = F(b) - F(a)
  \]
  Le calcul int\'{e}gral g\'{e}n\'{e}ralise la notion de primitive
  en donnant une \textbf{valeur num\'{e}rique} \`{a} l'int\'{e}grale.
\end{rappel}

%----------------------------------------------------------
\section{Int\'{e}grale d'une fonction continue}
%----------------------------------------------------------

\begin{definition}[Int\'{e}grale de $a$ \`{a} $b$]
  Soit $f$ une fonction \textbf{continue} sur $[a,b]$ et $F$
  une primitive de $f$ sur $[a,b]$.

  \medskip
  L'\textbf{int\'{e}grale de $f$ de $a$ \`{a} $b$} est le nombre r\'{e}el :
  \[
    \boxed{\int_a^b f(x)\,dx = \Bigl[F(x)\Bigr]_a^b = F(b) - F(a)}
  \]
\end{definition}

\begin{exemple}[R\'{e}solus : calculs d'int\'{e}grales]
  \begin{itemize}
    \item $\displaystyle\int_1^3 (2x+1)\,dx
      = \bigl[x^2+x\bigr]_1^3
      = (9+3) - (1+1) = 12 - 2 = 10$

    \medskip
    \item $\displaystyle\int_0^1 e^x\,dx
      = \bigl[e^x\bigr]_0^1
      = e^1 - e^0 = e - 1$

    \medskip
    \item $\displaystyle\int_1^e \frac{1}{x}\,dx
      = \bigl[\ln x\bigr]_1^e
      = \ln e - \ln 1 = 1 - 0 = 1$
  \end{itemize}
\end{exemple}

%----------------------------------------------------------
\section{Propri\'{e}t\'{e}s de l'int\'{e}grale}
%----------------------------------------------------------

\begin{propriete}[Propri\'{e}t\'{e}s fondamentales]
  \bigskip
  \begin{center}
  \renewcommand{\arraystretch}{2.2}
  \begin{tabular}{|p{4cm}|p{8cm}|}
    \hline
    \textbf{\color{navy}Propri\'{e}t\'{e}}
      & \textbf{\color{navy}Formule} \\
    \hline\hline
    Bornes \'{e}gales
      & $\displaystyle\int_a^a f(x)\,dx = 0$ \\
    \hline
    Inversion des bornes
      & $\displaystyle\int_b^a f(x)\,dx = -\int_a^b f(x)\,dx$ \\
    \hline
    Lin\'{e}arit\'{e}
      & $\displaystyle\int_a^b \bigl[f(x)+g(x)\bigr]dx
         = \int_a^b f(x)\,dx + \int_a^b g(x)\,dx$ \\
    \hline
    Homog\'{e}n\'{e}it\'{e}
      & $\displaystyle\int_a^b kf(x)\,dx = k\int_a^b f(x)\,dx$
        \quad ($k \in \R$) \\
    \hline
    Relation de Chasles
      & $\displaystyle\int_a^b f(x)\,dx
         = \int_a^c f(x)\,dx + \int_c^b f(x)\,dx$ \\
    \hline
  \end{tabular}
  \end{center}
\end{propriete}

\begin{propriete}[Int\'{e}grale et ordre]
  \begin{itemize}
    \item Si $f(x) \geq 0$ sur $[a,b]$, alors
          $\displaystyle\int_a^b f(x)\,dx \geq 0$
    \item Si $f(x) \geq g(x)$ sur $[a,b]$, alors
          $\displaystyle\int_a^b f(x)\,dx \geq \int_a^b g(x)\,dx$
  \end{itemize}
\end{propriete}

%----------------------------------------------------------
\section{Valeur moyenne}
%----------------------------------------------------------

\begin{definition}[Valeur moyenne d'une fonction]
  Soit $f$ continue sur $[a,b]$.
  La \textbf{valeur moyenne} de $f$ sur $[a,b]$ est :
  \[
    \mu = \frac{1}{b-a}\int_a^b f(x)\,dx
  \]
\end{definition}

\begin{exemple}[R\'{e}solu]
  Valeur moyenne de $f(x) = x^2$ sur $[0,3]$ :
  \[
    \mu = \frac{1}{3-0}\int_0^3 x^2\,dx
        = \frac{1}{3}\Bigl[\frac{x^3}{3}\Bigr]_0^3
        = \frac{1}{3} \times \frac{27}{3}
        = \frac{1}{3} \times 9 = 3
  \]
\end{exemple}

%----------------------------------------------------------
\section{Int\'{e}gration par parties}
%----------------------------------------------------------

\begin{theoreme}[Int\'{e}gration par parties (IPP)]
  Si $u$ et $v$ sont deux fonctions d\'{e}rivables sur $[a,b]$, alors :
  \[
    \boxed{\int_a^b u(x)\,v'(x)\,dx
    = \Bigl[u(x)\,v(x)\Bigr]_a^b - \int_a^b u'(x)\,v(x)\,dx}
  \]
\end{theoreme}

\begin{methode}[Choisir $u$ et $v'$ pour une IPP]
  \begin{enumerate}
    \item \textbf{Choisir} $u$ tel que $u'$ soit plus simple que $u$.
    \item \textbf{Choisir} $v'$ tel que $v$ soit facile \`{a} trouver.
    \item R\`{e}gle pratique : pr\'{e}f\'{e}rer $u = \ln$, $u = x^n$,
          et $v' = e^x$, $v' = \cos x$, $v' = \sin x$.
    \item Calculer $\bigl[u\,v\bigr]_a^b$ puis la nouvelle int\'{e}grale.
  \end{enumerate}
\end{methode}

\begin{exemple}[R\'{e}solu : IPP]
  Calculer $\displaystyle\int_0^1 x\,e^x\,dx$.

  \medskip
  On pose $u = x$ et $v' = e^x$,
  donc $u' = 1$ et $v = e^x$.

  \[
    \int_0^1 x\,e^x\,dx
    = \bigl[x\,e^x\bigr]_0^1 - \int_0^1 1 \cdot e^x\,dx
    = (1 \cdot e - 0) - \bigl[e^x\bigr]_0^1
    = e - (e - 1) = 1
  \]
\end{exemple}

\begin{exemple}[R\'{e}solu : IPP avec $\ln$]
  Calculer $\displaystyle\int_1^e \ln x\,dx$.

  \medskip
  On pose $u = \ln x$ et $v' = 1$,
  donc $u' = \dfrac{1}{x}$ et $v = x$.

  \[
    \int_1^e \ln x\,dx
    = \bigl[x\ln x\bigr]_1^e - \int_1^e x \cdot \frac{1}{x}\,dx
    = (e\ln e - 1\ln 1) - \int_1^e 1\,dx
    = e - 0 - \bigl[x\bigr]_1^e
    = e - (e - 1) = 1
  \]
\end{exemple}

%----------------------------------------------------------
\section{Calcul d'aires}
%----------------------------------------------------------

\begin{definition}[Unit\'{e} d'aire]
  Dans un rep\`{e}re orthogonal $(O,\vec{i},\vec{j})$,
  l'\textbf{unit\'{e} d'aire} (u.A.) est l'aire du rectangle bâti
  \`{a} partir des points $O$, $I$ et $J$ :
  \[
    1\,\text{u.A.} = \|\vec{i}\| \times \|\vec{j}\|
  \]
\end{definition}

\begin{propriete}[Calcul d'aires : cas g\'{e}n\'{e}ral]
  L'aire du domaine d\'{e}limit\'{e} par $\mathcal{C}_f$, l'axe des abscisses
  et les droites $x = a$, $x = b$ est :
  \[
    \mathcal{A} = \left(\int_a^b \bigl|f(x)\bigr|\,dx\right)\text{u.A.}
  \]
  L'aire du domaine d\'{e}limit\'{e} par $\mathcal{C}_f$, $\mathcal{C}_g$
  et les droites $x = a$, $x = b$ est :
  \[
    \mathcal{A} = \left(\int_a^b \bigl|f(x)-g(x)\bigr|\,dx\right)\text{u.A.}
  \]
\end{propriete}

\begin{propriete}[Cas particuliers selon le signe de $f$]
  \bigskip
  \begin{center}
  \renewcommand{\arraystretch}{2.2}
  \begin{tabular}{|p{6cm}|p{6.5cm}|}
    \hline
    \textbf{\color{navy}Situation}
      & \textbf{\color{navy}Aire} \\
    \hline\hline
    $f(x) \geq 0$ sur $[a,b]$
      & $\mathcal{A} = \left(\displaystyle\int_a^b f(x)\,dx\right)\text{u.A.}$ \\
    \hline
    $f(x) \leq 0$ sur $[a,b]$
      & $\mathcal{A} = \left(\displaystyle\int_a^b -f(x)\,dx\right)\text{u.A.}$ \\
    \hline
    $f \geq 0$ sur $[a,c]$ et $f \leq 0$ sur $[c,b]$
      & $\mathcal{A} = \left(\displaystyle\int_a^c f(x)\,dx
        + \int_c^b -f(x)\,dx\right)\text{u.A.}$ \\
    \hline
    $\mathcal{C}_f$ au-dessus de $\mathcal{C}_g$ sur $[a,b]$
      & $\mathcal{A} = \left(\displaystyle\int_a^b
        \bigl(f(x)-g(x)\bigr)dx\right)\text{u.A.}$ \\
    \hline
    $\mathcal{C}_f$ au-dessus de $\mathcal{C}_g$ sur $[a,c]$\\
    $\mathcal{C}_f$ au-dessous de $\mathcal{C}_g$ sur $[c,b]$
      & $\mathcal{A} = \left(\displaystyle\int_a^c(f-g)\,dx
        + \int_c^b(g-f)\,dx\right)\text{u.A.}$ \\
    \hline
  \end{tabular}
  \end{center}
\end{propriete}

\begin{methode}[Calculer une aire entre deux courbes]
  \begin{enumerate}
    \item Trouver les \textbf{points d'intersection} : r\'{e}soudre $f(x) = g(x)$.
    \item D\'{e}terminer \textbf{quelle courbe est au-dessus} sur chaque intervalle.
    \item Calculer $\displaystyle\int_a^b |f(x) - g(x)|\,dx$
          en d\'{e}coupant si n\'{e}cessaire.
    \item \textbf{Multiplier} par l'unit\'{e} d'aire.
  \end{enumerate}
\end{methode}

\begin{exemple}[R\'{e}solu : aire entre deux courbes]
  Aire entre $f(x) = x^2$ et $g(x) = x$ sur $[0,1]$.

  \medskip
  \textbf{\'{E}tape 1 :} $x^2 = x \Rightarrow x(x-1) = 0 \Rightarrow x = 0$ ou $x = 1$.

  \textbf{\'{E}tape 2 :} Sur $[0,1]$, $g(x) = x \geq x^2 = f(x)$.

  \textbf{\'{E}tape 3 :}
  \[
    \mathcal{A} = \int_0^1 (x - x^2)\,dx
    = \left[\frac{x^2}{2} - \frac{x^3}{3}\right]_0^1
    = \frac{1}{2} - \frac{1}{3} = \frac{1}{6}\;\text{u.A.}
  \]

  \begin{center}
  \begin{tikzpicture}
    \begin{axis}[
      axis lines  = center,
      xlabel      = {$x$},
      ylabel      = {$y$},
      xmin = -0.3, xmax = 1.5,
      ymin = -0.2, ymax = 1.3,
      xtick       = {0,0.5,1},
      ytick       = {0,0.5,1},
      grid        = major,
      grid style  = {gray!20},
      width = 9cm, height = 7cm,
    ]
      % Zone hachurée
      \addplot[fill=navy!15, draw=none, domain=0:1, samples=100]
        {x} \closedcycle;
      \addplot[fill=white, draw=none, domain=0:1, samples=100]
        {x^2} \closedcycle;

      % Courbe g(x) = x
      \addplot[forest, very thick, domain=-0.2:1.4]
        {x};
      \addlegendentry{$g(x)=x$}

      % Courbe f(x) = x²
      \addplot[burgundy, very thick, domain=-0.1:1.3, samples=100]
        {x^2};
      \addlegendentry{$f(x)=x^2$}

      % Points d'intersection
      \addplot[mark=*, mark size=3pt, black]
        coordinates {(0,0)(1,1)};

      % Annotation aire
      \node[navy, font=\small\bfseries] at (axis cs:0.5,0.4)
        {$\mathcal{A}=\frac{1}{6}$\,u.A.};
    \end{axis}
  \end{tikzpicture}
  \end{center}
\end{exemple}

%----------------------------------------------------------
\section{Calcul de volumes de r\'{e}volution}
%----------------------------------------------------------

\begin{definition}[Volume de r\'{e}volution]
  Le \textbf{volume} du solide engendr\'{e} par la rotation de
  $\mathcal{C}_f$ sur $[a,b]$ d'un tour complet autour de l'axe
  des abscisses est :
  \[
    \boxed{V = \pi \int_a^b \bigl[f(x)\bigr]^2\,dx \quad \text{u.v.}}
  \]
  o\`{u} u.v. d\'{e}signe l'unit\'{e} de volume.
\end{definition}

\begin{exemple}[R\'{e}solu : volume de r\'{e}volution]
  Volume engendr\'{e} par $f(x) = \sqrt{x}$ sur $[0,4]$
  autour de l'axe des abscisses :
  \[
    V = \pi\int_0^4 (\sqrt{x})^2\,dx
      = \pi\int_0^4 x\,dx
      = \pi\left[\frac{x^2}{2}\right]_0^4
      = \pi \times 8 = 8\pi\;\text{u.v.}
  \]
\end{exemple}

\newpage
%----------------------------------------------------------
\section*{Exercices du chapitre}
\addcontentsline{toc}{section}{Exercices : Calcul int\'{e}gral}
%----------------------------------------------------------

\begin{exercice}[1 : Calculs d'int\'{e}grales]
  Calculer les int\'{e}grales suivantes :
  \begin{multicols}{2}
  \begin{enumerate}
    \item $\displaystyle\int_0^2 (3x^2-2x+1)\,dx$
    \item $\displaystyle\int_1^4 \frac{1}{\sqrt{x}}\,dx$
    \item $\displaystyle\int_0^{\ln 2} e^x\,dx$
    \item $\displaystyle\int_1^e \frac{2}{x}\,dx$
    \item $\displaystyle\int_0^{\pi} \cos x\,dx$
    \item $\displaystyle\int_0^1 \frac{x}{x^2+1}\,dx$
  \end{enumerate}
  \end{multicols}
\end{exercice}

\begin{exercice}[2 : Propri\'{e}t\'{e}s]
  \begin{enumerate}
    \item Sachant que $\displaystyle\int_0^3 f(x)\,dx = 7$ et
          $\displaystyle\int_0^3 g(x)\,dx = 2$, calculer
          $\displaystyle\int_0^3 [3f(x) - 2g(x)]\,dx$.
    \item Sachant que $\displaystyle\int_0^5 f(x)\,dx = 10$ et
          $\displaystyle\int_0^3 f(x)\,dx = 4$, calculer
          $\displaystyle\int_3^5 f(x)\,dx$.
  \end{enumerate}
\end{exercice}

\begin{exercice}[3 : Valeur moyenne]
  Calculer la valeur moyenne de :
  \begin{enumerate}
    \item $f(x) = x^2 + 1$ sur $[0,3]$
    \item $f(x) = e^x$ sur $[0,2]$
    \item $f(x) = \dfrac{1}{x}$ sur $[1,e]$
  \end{enumerate}
\end{exercice}

\begin{exercice}[4 : Int\'{e}gration par parties]
  Calculer par une int\'{e}gration par parties :
  \begin{enumerate}
    \item $\displaystyle\int_0^1 x\,e^{2x}\,dx$
    \item $\displaystyle\int_1^e x\ln x\,dx$
    \item $\displaystyle\int_0^{\pi} x\sin x\,dx$
    \item $\displaystyle\int_0^1 (x+1)e^x\,dx$
  \end{enumerate}
\end{exercice}

\begin{exercice}[5 : Calcul d'aires]
  \begin{enumerate}
    \item Calculer l'aire du domaine d\'{e}limit\'{e} par $\mathcal{C}_f$,
          l'axe des abscisses et les droites $x=0$, $x=2$
          pour $f(x) = x^2 - x$.
    \item Calculer l'aire entre $f(x) = x^2$ et $g(x) = \sqrt{x}$
          sur $[0,1]$.
    \item Calculer l'aire entre $f(x) = e^x$ et $g(x) = x + 1$
          sur $[0,2]$.
  \end{enumerate}
\end{exercice}

\begin{exercice}[6 : Volume de r\'{e}volution]
  \begin{enumerate}
    \item Calculer le volume engendr\'{e} par $f(x) = x$ sur $[0,3]$
          autour de l'axe des abscisses.
          \textit{(V\'{e}rifier : on doit trouver le volume d'un c\^{o}ne.)}
    \item Calculer le volume engendr\'{e} par $f(x) = e^x$ sur $[0,1]$.
  \end{enumerate}
\end{exercice}
% !TEX root = ../main.tex
%==========================================================
%  CHAPITRE 9 — PROBABILITÉS
%==========================================================

\chapter{Probabilit\'{e}s}

%----------------------------------------------------------
\section{Probabilit\'{e}s d'un ensemble fini}
%----------------------------------------------------------

\begin{definition}[Probabilit\'{e} d'un \'{e}v\'{e}nement]
La \textbf{probabilit\'{e}} d'un \'{e}v\'{e}nement $A$ est la somme des probabilit\'{e}s
des \'{e}v\'{e}nements \'{e}l\'{e}mentaires qui le composent. On la note $p(A)$.
\end{definition}

\begin{propriete}[Propri\'{e}t\'{e}s fondamentales]
Soit $\Omega$ l'univers associ\'{e} \`{a} une exp\'{e}rience al\'{e}atoire.

\medskip
\begin{center}
\renewcommand{\arraystretch}{1.5}
\begin{tabular}{|c|c|}
  \hline
  \rowcolor{forestlight}
  \textbf{\color{forest}L'\'{e}v\'{e}nement} & \textbf{\color{forest}Probabilit\'{e}} \\
  \hline
  $A$ & $0 \leq p(A) \leq 1$ \\
  \hline
  $\bar{A}$ & $p(\bar{A}) = 1 - p(A)$ \\
  \hline
  $A \cup B$ & $p(A \cup B) = p(A) + p(B) - p(A \cap B)$ \\
  \hline
  $A \cup B$ \quad ($A$ et $B$ incompatibles) & $p(A \cup B) = p(A) + p(B)$ \\
  \hline
\end{tabular}
\end{center}
\end{propriete}

\begin{theoreme}[Equiprobabilit\'{e}]
S'il y a \textbf{\'{e}quiprobabilit\'{e}}, alors pour tout \'{e}v\'{e}nement $A$ de $\Omega$ :
\[
  p(A) = \frac{\mathrm{card}(A)}{\mathrm{card}(\Omega)}
       = \frac{\text{nombre des cas favorables}}{\text{nombre des cas possibles}}
\]
\end{theoreme}

%----------------------------------------------------------
\section{Variable al\'{e}atoire et loi de probabilit\'{e}}
%----------------------------------------------------------

\begin{definition}[Loi de probabilit\'{e} d'une variable al\'{e}atoire]
Soit $\Omega$ l'univers associ\'{e} \`{a} une exp\'{e}rience al\'{e}atoire.
Pour d\'{e}finir la loi de probabilit\'{e} de la variable $X$ sur $\Omega$ :
\begin{enumerate}[label=\textbf{\arabic*.}, itemsep=2pt]
  \item On d\'{e}termine $X(\Omega) = \{x_1;\,x_2;\,x_3;\,\ldots;\,x_n\}$
        l'ensemble des valeurs prises par $X$.
  \item On calcule pour chaque valeur $x_i$ sa probabilit\'{e}
        $p_i = p(X = x_i)$ avec $i \in \{1;\,2;\,\ldots;\,n\}$.
  \item On r\'{e}sume la loi de probabilit\'{e} de $X$ par le tableau :
\end{enumerate}

\medskip
\begin{center}
\renewcommand{\arraystretch}{1.5}
\begin{tabular}{|c|c|c|c|c|c|}
  \hline
  \rowcolor{navylight}
  $x_i$ & $x_1$ & $x_2$ & $x_3$ & $\cdots$ & $x_n$ \\
  \hline
  $p(X = x_i)$ & $p_1$ & $p_2$ & $p_3$ & $\cdots$ & $p_n$ \\
  \hline
\end{tabular}
\end{center}

\medskip
\noindent Avec la condition : $p_1 + p_2 + \cdots + p_n = 1$.
\end{definition}

%----------------------------------------------------------
\section{Probabilit\'{e} conditionnelle et ind\'{e}pendance}
%----------------------------------------------------------

\begin{definition}[Probabilit\'{e} conditionnelle]
Soit $A$ et $B$ deux \'{e}v\'{e}nements li\'{e}s \`{a} une m\^{e}me exp\'{e}rience al\'{e}atoire
tels que $p(A) \neq 0$.\\
La \textbf{probabilit\'{e} de l'\'{e}v\'{e}nement $B$ sachant que $A$ est r\'{e}alis\'{e}}
est le nombre :
\[
  p_A(B) = p\!\left(\frac{B}{A}\right) = \frac{p(A \cap B)}{p(A)}
\]
\end{definition}

\begin{theoreme}[Ev\'{e}nements ind\'{e}pendants]
Soit $A$ et $B$ deux \'{e}v\'{e}nements li\'{e}s \`{a} une m\^{e}me exp\'{e}rience al\'{e}atoire.\\
$A$ et $B$ sont \textbf{ind\'{e}pendants} $\Longleftrightarrow$
\[
  p(A \cap B) = p(A) \times p(B)
\]
\end{theoreme}

\begin{remarque}
Si $A$ et $B$ sont ind\'{e}pendants, alors $p_A(B) = p(B)$ :
la r\'{e}alisation de $A$ n'influe pas sur la probabilit\'{e} de $B$.
\end{remarque}

%----------------------------------------------------------
\section{Esp\'{e}rance, variance et \'{e}cart-type}
%----------------------------------------------------------

\begin{definition}[Esp\'{e}rance, variance, \'{e}cart-type]
Soit $X$ une variable al\'{e}atoire dont la loi de probabilit\'{e} est :

\medskip
\begin{center}
\renewcommand{\arraystretch}{1.5}
\begin{tabular}{|c|c|c|c|c|c|}
  \hline
  \rowcolor{navylight}
  $x_i$ & $x_1$ & $x_2$ & $x_3$ & $\cdots$ & $x_n$ \\
  \hline
  $p(X = x_i)$ & $p_1$ & $p_2$ & $p_3$ & $\cdots$ & $p_n$ \\
  \hline
\end{tabular}
\end{center}

\medskip
\begin{center}
\renewcommand{\arraystretch}{1.8}
\begin{tabular}{|l|l|}
  \hline
  \rowcolor{navylight}
  \textbf{\color{navy}Grandeur} & \textbf{\color{navy}Formule} \\
  \hline
  Esp\'{e}rance math\'{e}matique de $X$
    & $E(X) = x_1 p_1 + x_2 p_2 + \cdots + x_n p_n$ \\
  \hline
  Variance de $X$
    & $V(X) = E(X^2) - \bigl[E(X)\bigr]^2$ \\
  \hline
  Ecart-type de $X$
    & $\sigma(X) = \sqrt{V(X)}$ \\
  \hline
\end{tabular}
\end{center}
\end{definition}

\begin{exemple}
Soit $X$ une variable al\'{e}atoire de loi :

\begin{center}
\renewcommand{\arraystretch}{1.4}
\begin{tabular}{|c|c|c|c|}
  \hline
  \rowcolor{navylight}
  $k$ & $-1$ & $0$ & $2$ \\
  \hline
  $p(X=k)$ & $\dfrac{1}{6}$ & $\dfrac{1}{3}$ & $\dfrac{1}{2}$ \\
  \hline
\end{tabular}
\end{center}

\medskip
$E(X) = (-1)\cdot\dfrac{1}{6} + 0\cdot\dfrac{1}{3} + 2\cdot\dfrac{1}{2}
       = -\dfrac{1}{6} + 1 = \dfrac{5}{6}$

$E(X^2) = 1\cdot\dfrac{1}{6} + 0\cdot\dfrac{1}{3} + 4\cdot\dfrac{1}{2}
         = \dfrac{1}{6} + 2 = \dfrac{13}{6}$

$V(X) = \dfrac{13}{6} - \left(\dfrac{5}{6}\right)^2
       = \dfrac{13}{6} - \dfrac{25}{36} = \dfrac{53}{36}$
\end{exemple}

%----------------------------------------------------------
\section{Epreuve r\'{e}p\'{e}t\'{e}e et loi binomiale}
%----------------------------------------------------------

\begin{theoreme}[Epreuve r\'{e}p\'{e}t\'{e}e]
Soit $p$ la probabilit\'{e} d'un \'{e}v\'{e}nement $A$ lors d'une exp\'{e}rience al\'{e}atoire.
Si on r\'{e}p\`{e}te $n$ fois l'\'{e}preuve dans des conditions identiques,
alors la probabilit\'{e} de r\'{e}alisation de $A$ \textbf{exactement $k$ fois}
durant les $n$ \'{e}preuves est :
\[
  C_n^k \cdot p^k \cdot (1-p)^{n-k} \qquad (k \leq n)
\]
\end{theoreme}

\begin{definition}[Loi binomiale]
Soit $p$ la probabilit\'{e} d'un \'{e}v\'{e}nement $A$ lors d'une exp\'{e}rience al\'{e}atoire.
On r\'{e}p\`{e}te $n$ fois l'\'{e}preuve dans des conditions identiques.
Si la variable al\'{e}atoire $X$ est \'{e}gale au nombre de r\'{e}alisations de $A$
durant les $n$ \'{e}preuves, alors :
\[
  \forall k \in \{0;\,1;\,2;\,\ldots;\,n\} \qquad
  p(X = k) = C_n^k \cdot p^k \cdot (1-p)^{n-k}
\]
On dit que $X$ suit une \textbf{loi binomiale} de param\`{e}tres $n$ et $p$,
not\'{e}e $X \sim \mathcal{B}(n,p)$.
\end{definition}

\begin{propriete}[Esp\'{e}rance et variance de la loi binomiale]
Si $X \sim \mathcal{B}(n,p)$, alors :
\[
  E(X) = n \cdot p \qquad \text{et} \qquad V(X) = n \cdot p \cdot (1-p)
\]
\end{propriete}

\begin{methode}[Reconna\^{i}tre et utiliser la loi binomiale]
\begin{enumerate}[label=\textbf{\arabic*.}, itemsep=3pt]
  \item V\'{e}rifier que l'exp\'{e}rience est r\'{e}p\'{e}t\'{e}e $n$ fois dans des
        \textbf{conditions identiques et ind\'{e}pendantes}.
  \item Identifier l'\'{e}v\'{e}nement $A$ de probabilit\'{e} $p$ (\emph{succ\`{e}s}).
  \item Identifier $X$ = nombre de succ\`{e}s $\Rightarrow X \sim \mathcal{B}(n,p)$.
  \item Calculer $p(X = k) = C_n^k \cdot p^k \cdot (1-p)^{n-k}$.
\end{enumerate}
\end{methode}

\begin{exemple}
On lance un d\'{e} \'{e}quilibr\'{e} $5$ fois. Soit $X$ le nombre de fois o\`{u} on
obtient un $6$. Alors $X \sim \mathcal{B}\!\left(5,\dfrac{1}{6}\right)$.

\medskip
$p(X = 2) = C_5^2 \cdot \left(\dfrac{1}{6}\right)^2 \cdot \left(\dfrac{5}{6}\right)^3
           = 10 \cdot \dfrac{1}{36} \cdot \dfrac{125}{216}
           = \dfrac{1250}{7776} \approx 0{,}161$

\medskip
$E(X) = 5 \times \dfrac{1}{6} = \dfrac{5}{6}$ \qquad
$V(X) = 5 \times \dfrac{1}{6} \times \dfrac{5}{6} = \dfrac{25}{36}$
\end{exemple}

%----------------------------------------------------------
\section{Exercices d'entra\^{i}nement}
%----------------------------------------------------------

\begin{exercice}[1]
Une urne contient 4 boules rouges et 6 boules bleues.
On tire simultan\'{e}ment 3 boules.

\begin{enumerate}[label=\textbf{\arabic*.}, itemsep=4pt]
  \item Calculer la probabilit\'{e} d'obtenir exactement 2 boules rouges.
  \item Calculer la probabilit\'{e} d'obtenir au moins une boule rouge.
\end{enumerate}
\end{exercice}

\begin{exercice}[2]
Soit $X$ une variable al\'{e}atoire dont la loi est :

\begin{center}
\renewcommand{\arraystretch}{1.4}
\begin{tabular}{|c|c|c|c|c|}
  \hline
  \rowcolor{navylight}
  $k$ & $0$ & $1$ & $2$ & $3$ \\
  \hline
  $p(X=k)$ & $0{,}1$ & $a$ & $2a$ & $0{,}2$ \\
  \hline
\end{tabular}
\end{center}

\begin{enumerate}[label=\textbf{\arabic*.}, itemsep=4pt]
  \item D\'{e}terminer $a$.
  \item Calculer $E(X)$, $V(X)$ et $\sigma(X)$.
\end{enumerate}
\end{exercice}

\begin{exercice}[3]
Un contr\^{o}le de qualit\'{e} montre que $5\,\%$ des pi\`{e}ces produites par
une machine sont d\'{e}fectueuses. On pr\'{e}l\`{e}ve au hasard $10$ pi\`{e}ces.
Soit $X$ le nombre de pi\`{e}ces d\'{e}fectueuses.

\begin{enumerate}[label=\textbf{\arabic*.}, itemsep=4pt]
  \item Quelle est la loi suivie par $X$ ? Pr\'{e}ciser ses param\`{e}tres.
  \item Calculer $p(X = 0)$ et $p(X \geq 1)$.
  \item Calculer $E(X)$ et $V(X)$.
\end{enumerate}
\end{exercice}


%----------------------------------------------------------
%  COMMANDES SPÉCIALES EXAMENS (lignes pointillées + boîte)
%----------------------------------------------------------
\newcommand{\lignerep}{%
  \par\noindent\makebox[\linewidth]{\dotfill}\par
}
\newtcolorbox{exerciceBox}[2]{
  colback=navylight, colframe=navy,
  boxrule=1pt, arc=3pt,
  left=8pt, right=8pt, top=5pt, bottom=5pt,
  coltitle=white, colbacktitle=navy,
  fonttitle=\bfseries\large,
  title={\textit{Exercice #1 :} \hfill (#2 points)},
  toptitle=3pt, bottomtitle=3pt
}

%----------------------------------------------------------
%  EXAMENS BLANCS
%----------------------------------------------------------
 \cleardoublepage
 \part*{Examens Blancs}
 \addcontentsline{toc}{part}{Examens Blancs}

 % !TEX root = ../main.tex
%==========================================================
%  EXAMEN BLANC N°1 — VERSION COMPACTE (LIVRE)
%==========================================================

\cleardoublepage
\addcontentsline{toc}{chapter}{Examen Blanc N\textsuperscript{o}1}
\thispagestyle{empty}

%----------------------------------------------------------
%  TITRE
%----------------------------------------------------------
\begin{center}
  {\LARGE\bfseries\color{navy} Examen Blanc N\textsuperscript{o}1}\\[0.3cm]
  {\large\bfseries Math\'{e}matiques \textemdash{} Bac Sciences \'{E}conomiques \textemdash{} Session 2026}
\end{center}

\smallskip
{\small Version compacte professionnelle adapt\'{e}e pour un livre de pr\'{e}paration au baccalaur\'{e}at.}
\smallskip

%----------------------------------------------------------
%  TABLEAU DES COMPOSANTES
%----------------------------------------------------------
\begin{center}
\renewcommand{\arraystretch}{1.3}
\begin{tabular}{|c|l|c|}
  \hline
  \rowcolor{navylight}
  \textbf{\color{navy}Exercice} & \multicolumn{1}{c|}{\textbf{\color{navy}Th\`{e}me}} & \textbf{\color{navy}Points} \\
  \hline
  1 & Suites num\'{e}riques & 4 \\
  \hline
  2 & Probabilit\'{e}s & 1,5 \\
  \hline
  3 & Probabilit\'{e}s & 2,5 \\
  \hline
  4 & \'{E}tude graphique & 4,75 \\
  \hline
  5 & Fonction et int\'{e}grale & 7,25 \\
  \hline
\end{tabular}
\end{center}

\bigskip

%==========================================================
%  EXERCICE 1 — SUITES NUMÉRIQUES  (4 pts)
%==========================================================
\noindent\textbf{Exercice 1 :}\\
Soit la suite d\'{e}finie par :
$u_0 = 1$ \quad et \quad $u_{n+1} = 6 - \dfrac{7}{u_n + 2}$ \quad pour tout $n \in \N$.

\begin{enumerate}[leftmargin=1.5em, label=\textbf{\arabic*)}, itemsep=2pt, topsep=4pt]
  \item Calculer $u_1$ et $u_2$.
  \item Montrer par r\'{e}currence que $1 \leq u_n \leq 5$ pour tout $n \in \N$.
  \item
    \begin{enumerate}[label=\textbf{\alph*)}, itemsep=1pt, topsep=1pt]
      \item V\'{e}rifier que $u_{n+1} = \dfrac{6u_n + 5}{u_n + 2}$ pour tout $n \in \N$.
      \item Montrer que $u_{n+1} - u_n = \dfrac{-(u_n-5)(1+u_n)}{u_n+2}$,
            puis d\'{e}duire que $(u_n)$ est croissante.
      \item En d\'{e}duire que $(u_n)$ est convergente.
    \end{enumerate}
  \item On consid\`{e}re $(v_n)$ d\'{e}finie par $v_n = \dfrac{u_n - 5}{u_n + 1}$.
    \begin{enumerate}[label=\textbf{\alph*)}, itemsep=1pt, topsep=1pt]
      \item Montrer que $(v_n)$ est g\'{e}om\'{e}trique de raison $\dfrac{1}{7}$.
      \item D\'{e}terminer $v_n$ en fonction de $n$, puis $u_n$ en fonction de $n$.
      \item Calculer $\displaystyle\lim_{n \to +\infty} u_n$.
      \item Calculer $S_n = v_0 + v_1 + \cdots + v_n$ en fonction de $n$.
    \end{enumerate}
\end{enumerate}

\medskip

%==========================================================
%  EXERCICE 2 — PROBABILITÉS  (1,5 pts)
%==========================================================
\noindent\textbf{Exercice 2 :}\\
Soit $X$ une variable al\'{e}atoire telle que :
$p(X=-1)=a$ ; $p(X=1)=2a$ ; $p(X=2)=3a$.

\begin{enumerate}[leftmargin=1.5em, label=\textbf{\arabic*.}, itemsep=2pt, topsep=4pt]
  \item D\'{e}terminer $a$.
  \item Calculer $E(X)$ et $V(X)$.
\end{enumerate}

\medskip

%==========================================================
%  EXERCICE 3 — PROBABILITÉS  (2,5 pts)
%==========================================================
\noindent\textbf{Exercice 3 :}\\
Une urne contient 5 boules noires num\'{e}rot\'{e}es $1;1;2;2;2$
et 4 boules blanches num\'{e}rot\'{e}es $1;1;2;2$.
On tire simultan\'{e}ment 3 boules. On pose :
\begin{itemize}[itemsep=0pt, topsep=2pt, leftmargin=2em]
  \item $A$ : \guillemotleft{} les 3 boules sont de m\^{e}me couleur \guillemotright{}
  \item $B$ : \guillemotleft{} les 3 boules portent le m\^{e}me nombre \guillemotright{}
\end{itemize}

\begin{enumerate}[leftmargin=1.5em, label=\textbf{\arabic*.}, itemsep=2pt, topsep=4pt]
  \item Calculer $p(A)$, $p(B)$ puis montrer que $p(A \cap B) = \dfrac{1}{84}$.
  \item Les \'{e}v\'{e}nements $A$ et $B$ sont-ils ind\'{e}pendants ?
  \item Sachant que les boules sont de m\^{e}me couleur,
        calculer la probabilit\'{e} qu'elles portent le m\^{e}me nombre.
\end{enumerate}

\medskip

%==========================================================
%  EXERCICE 4 — ÉTUDE GRAPHIQUE  (4,75 pts)
%==========================================================
\noindent\textbf{Exercice 4 :}\\
On consid\`{e}re la fonction $f$ d\'{e}finie par sa courbe repr\'{e}sentative $(\mathcal{C})$
dans un rep\`{e}re orthonorm\'{e}, avec une asymptote verticale $x = -1$
et une asymptote oblique $y = \dfrac{3}{5}\,x$.\\
R\'{e}pondre graphiquement aux questions suivantes :

\begin{enumerate}[leftmargin=1.5em, label=\textbf{\arabic*.}, itemsep=2pt, topsep=4pt]
  \item D\'{e}terminer le domaine de d\'{e}finition de $f$.
  \item D\'{e}terminer le signe de $f$.
  \item D\'{e}terminer $\displaystyle\lim_{x \to +\infty} f(x)$ et
        $\displaystyle\lim_{x \to -1^-} f(x)$,
        puis les branches infinies de $(\mathcal{C})$.
\end{enumerate}

\medskip

%==========================================================
%  EXERCICE 5 — ÉTUDE DE FONCTION  (7,25 pts)
%==========================================================
\noindent\textbf{Exercice 5 :}\\
Soit $f$ la fonction d\'{e}finie sur $\R$ par $f(x) = (1-x)^2\,e^x$.
Soit $(\mathcal{C}_f)$ sa courbe dans un rep\`{e}re orthonorm\'{e}
avec $\|\vec{i}\| = 1\,\text{cm}$.

\begin{enumerate}[leftmargin=1.5em, label=\textbf{\arabic*.}, itemsep=2pt, topsep=4pt]
  \item Calculer $\displaystyle\lim_{x \to +\infty} f(x)$ et
        $\displaystyle\lim_{x \to +\infty} \dfrac{f(x)}{x}$,
        puis donner une interpr\'{e}tation g\'{e}om\'{e}trique.
  \item V\'{e}rifier que $f(x) = x^2 e^x\!\left(1 - \dfrac{2}{x} + \dfrac{1}{x^2}\right)$,
        calculer $\displaystyle\lim_{x \to -\infty} f(x)$ et donner une interpr\'{e}tation
        g\'{e}om\'{e}trique.
  \item Montrer que $f'(x) = (x-1)(x+1)\,e^x$, puis dresser le tableau de variations de $f$.
  \item D\'{e}terminer l'\'{e}quation de la tangente $(T)$ \`{a} $(\mathcal{C}_f)$
        au point d'abscisse $0$.
  \item Montrer que $F : x \mapsto e^x(x^2 - 4x + 5)$ est une primitive de $f$,
        puis calculer l'aire (en cm\textsuperscript{2}) du domaine d\'{e}limit\'{e}
        par $(\mathcal{C}_f)$, la droite $y = 1$ et les droites
        $x = -\dfrac{5}{2}$, $x = \dfrac{3}{2}$.
\end{enumerate}

\begin{center}
  \color{charcoal}\rule{5cm}{0.4pt}\\[3pt]
  \textit{--- Fin du sujet ---}\\
  \rule{5cm}{0.4pt}
\end{center}

 % !TEX root = ../main.tex
%==========================================================
%  EXAMEN BLANC N°2 — VERSION COMPACTE (LIVRE)
%==========================================================

\cleardoublepage
\addcontentsline{toc}{chapter}{Examen Blanc N\textsuperscript{o}2}
\thispagestyle{empty}

%----------------------------------------------------------
%  TITRE
%----------------------------------------------------------
\begin{center}
  {\LARGE\bfseries\color{navy} Examen Blanc N\textsuperscript{o}2}\\[0.3cm]
  {\large\bfseries Math\'{e}matiques \textemdash{} Bac Sciences \'{E}conomiques \textemdash{} Session 2026}
\end{center}

\smallskip
{\small Version compacte professionnelle adapt\'{e}e pour un livre de pr\'{e}paration au baccalaur\'{e}at.}
\smallskip

%----------------------------------------------------------
%  TABLEAU DES COMPOSANTES
%----------------------------------------------------------
\begin{center}
\renewcommand{\arraystretch}{1.3}
\begin{tabular}{|c|l|c|}
  \hline
  \rowcolor{navylight}
  \textbf{\color{navy}Exercice} & \multicolumn{1}{c|}{\textbf{\color{navy}Th\`{e}me}} & \textbf{\color{navy}Points} \\
  \hline
  1 & Suites num\'{e}riques & 4 \\
  \hline
  2 & Probabilit\'{e}s & 1,5 \\
  \hline
  3 & Probabilit\'{e}s & 2,5 \\
  \hline
  4 & \'{E}tude graphique & 4,75 \\
  \hline
  5 & Fonction et int\'{e}grale & 7,25 \\
  \hline
\end{tabular}
\end{center}

\bigskip

%==========================================================
%  EXERCICE 1 — SUITES NUMÉRIQUES  (4 pts)
%==========================================================
\noindent\textbf{Exercice 1 :}\\
Soit $(u_n)$ une suite num\'{e}rique d\'{e}finie par :
\[
  u_0 = 1 \qquad \text{et} \qquad
  u_{n+1} = \frac{3u_n + 2}{u_n + 2}
  \quad \text{pour tout } n \in \N
\]

\begin{enumerate}[leftmargin=1.5em, label=\textbf{\arabic*)}, itemsep=2pt, topsep=4pt]
  \item Calculer $u_1$ et $u_2$.
  \item Montrer par r\'{e}currence que $1 \leq u_n \leq 2$ pour tout $n \in \N$.
  \item
    \begin{enumerate}[label=\textbf{\alph*)}, itemsep=1pt, topsep=1pt]
      \item V\'{e}rifier que $u_{n+1} - 2 = \dfrac{u_n - 2}{u_n + 2}$ pour tout $n \in \N$.
      \item Montrer que $u_{n+1} - u_n = \dfrac{-(u_n - 2)(u_n + 1)}{u_n + 2}$,
            puis d\'{e}duire que $(u_n)$ est croissante.
      \item En d\'{e}duire que $(u_n)$ est convergente.
    \end{enumerate}
  \item On consid\`{e}re $(v_n)$ d\'{e}finie par $v_n = \dfrac{u_n - 2}{u_n + 1}$ pour tout $n \in \N$.
    \begin{enumerate}[label=\textbf{\alph*)}, itemsep=1pt, topsep=1pt]
      \item Montrer que $(v_n)$ est une suite g\'{e}om\'{e}trique de raison $\dfrac{1}{4}$.
      \item D\'{e}terminer $v_n$ en fonction de $n$, puis $u_n$ en fonction de $n$.
      \item Calculer $\displaystyle\lim_{n \to +\infty} u_n$.
      \item Calculer $S_n = v_0 + v_1 + \cdots + v_n$ en fonction de $n$.
    \end{enumerate}
\end{enumerate}

\medskip

%==========================================================
%  EXERCICE 2 — PROBABILITÉS  (1,5 pts)
%==========================================================
\noindent\textbf{Exercice 2 :}\\
Soit $X$ une variable al\'{e}atoire dont la loi de probabilit\'{e} est :

\medskip
\begin{center}
\renewcommand{\arraystretch}{1.4}
\begin{tabular}{|c|c|c|c|}
  \hline
  \rowcolor{navylight}
  \textbf{\color{navy}$k$} & $-1$ & $0$ & $2$ \\
  \hline
  \textbf{\color{navy}$p(X=k)$} & $2a$ & $a$ & $3a$ \\
  \hline
\end{tabular}
\end{center}
\medskip

\begin{enumerate}[leftmargin=1.5em, label=\textbf{\arabic*.}, itemsep=2pt, topsep=4pt]
  \item D\'{e}terminer $a$.
  \item Calculer $E(X)$ et $V(X)$.
\end{enumerate}

\medskip

%==========================================================
%  EXERCICE 3 — PROBABILITÉS  (2,5 pts)
%==========================================================
\noindent\textbf{Exercice 3 :}\\
Une urne contient 4 boules rouges num\'{e}rot\'{e}es $1;1;2;2$
et 5 boules vertes num\'{e}rot\'{e}es $1;2;2;2;3$.
On tire simultan\'{e}ment 3 boules. On pose :
\begin{itemize}[itemsep=0pt, topsep=2pt, leftmargin=2em]
  \item $A$ : \guillemotleft{} les 3 boules tir\'{e}es sont de m\^{e}me couleur \guillemotright{}
  \item $B$ : \guillemotleft{} les 3 boules tir\'{e}es portent le m\^{e}me num\'{e}ro \guillemotright{}
\end{itemize}

\begin{enumerate}[leftmargin=1.5em, label=\textbf{\arabic*.}, itemsep=2pt, topsep=4pt]
  \item Calculer $p(A)$, $p(B)$ puis montrer que $p(A \cap B) = \dfrac{1}{84}$.
  \item Les \'{e}v\'{e}nements $A$ et $B$ sont-ils ind\'{e}pendants ?
  \item Sachant que les boules tir\'{e}es sont de m\^{e}me couleur,
        calculer la probabilit\'{e} qu'elles portent le m\^{e}me num\'{e}ro.
\end{enumerate}

\medskip

%==========================================================
%  EXERCICE 4 — ÉTUDE GRAPHIQUE  (4,75 pts)
%==========================================================
\noindent\textbf{Exercice 4 :}\\
On consid\`{e}re la fonction $f$ d\'{e}finie par sa courbe repr\'{e}sentative $(\mathcal{C})$
dans un rep\`{e}re orthonorm\'{e}, avec une asymptote verticale $x = 2$
et une asymptote oblique $y = \dfrac{2x+3}{3}$.\\
R\'{e}pondre graphiquement aux questions suivantes :

\begin{enumerate}[leftmargin=1.5em, label=\textbf{\arabic*.}, itemsep=2pt, topsep=4pt]
  \item D\'{e}terminer le domaine de d\'{e}finition de $f$.
  \item D\'{e}terminer le signe de $f$ sur son domaine.
  \item
    \begin{enumerate}[label=\textbf{\alph*.}, itemsep=1pt, topsep=1pt]
      \item D\'{e}terminer $\displaystyle\lim_{x \to +\infty} f(x)$
            et $\displaystyle\lim_{x \to 2^+} f(x)$.
      \item D\'{e}terminer les branches infinies de $(\mathcal{C})$.
    \end{enumerate}
\end{enumerate}

\newpage

%==========================================================
%  EXERCICE 5 — ÉTUDE DE FONCTION  (7,25 pts)
%==========================================================
\noindent\textbf{Exercice 5 :}\\
Soit $g$ la fonction d\'{e}finie sur $\R$ par $g(x) = (x+2)^2\,e^{-x}$.
Soit $(\mathcal{C}_g)$ sa courbe repr\'{e}sentative dans un rep\`{e}re
orthonorm\'{e} $(O;\,\vec{i};\,\vec{j})$, avec $\|\vec{i}\| = 1\,\text{cm}$.

\begin{enumerate}[leftmargin=1.5em, label=\textbf{\arabic*.}, itemsep=2pt, topsep=4pt]
  \item Calculer $\displaystyle\lim_{x \to +\infty} g(x)$ et donner une
        interpr\'{e}tation g\'{e}om\'{e}trique.
  \item V\'{e}rifier que $g(x) = x^2\,e^{-x}\!\left(1 + \dfrac{2}{x}\right)^2$
        pour tout $x \neq 0$, puis calculer $\displaystyle\lim_{x \to -\infty} g(x)$
        et $\displaystyle\lim_{x \to -\infty} \dfrac{g(x)}{x^2}$.
        Donner une interpr\'{e}tation g\'{e}om\'{e}trique.
  \item
    \begin{enumerate}[label=\textbf{\alph*.}, itemsep=1pt, topsep=1pt]
      \item Montrer que $\forall x \in \R$ :
            $g'(x) = -x(x+2)\,e^{-x}$.
      \item Dresser le tableau de variations de $g$ sur $\R$.
    \end{enumerate}
  \item D\'{e}terminer l'\'{e}quation de la tangente $(T)$ \`{a}
        $(\mathcal{C}_g)$ au point d'abscisse $0$.
  \item Montrer que la fonction $G : x \mapsto -(x^2 + 6x + 10)\,e^{-x}$
        est une primitive de $g$ sur $\R$.
  \item Calculer l'aire (en cm\textsuperscript{2}) du domaine plan d\'{e}limit\'{e}
        par $(\mathcal{C}_g)$, l'axe des abscisses et les droites
        $x = -2$ et $x = 1$.
  \item Montrer que $g(x) \leq 4$ pour tout $x \in [-2\,;\,+\infty[$.
        En d\'{e}duire la position de $(\mathcal{C}_g)$ par rapport
        \`{a} la tangente $(T)$ sur cet intervalle.
\end{enumerate}

\begin{center}
  \color{charcoal}\rule{5cm}{0.4pt}\\[3pt]
  \textit{--- Fin du sujet ---}\\
  \rule{5cm}{0.4pt}
\end{center}

 % !TEX root = ../main.tex
%==========================================================
%  EXAMEN BLANC N°3 — VERSION COMPACTE (LIVRE)
%==========================================================

\cleardoublepage
\addcontentsline{toc}{chapter}{Examen Blanc N\textsuperscript{o}3}
\thispagestyle{empty}

%----------------------------------------------------------
%  TITRE
%----------------------------------------------------------
\begin{center}
  {\LARGE\bfseries\color{navy} Examen Blanc N\textsuperscript{o}3}\\[0.3cm]
  {\large\bfseries Math\'{e}matiques \textemdash{} Bac Sciences \'{E}conomiques \textemdash{} Session 2026}
\end{center}

\smallskip
{\small Version compacte professionnelle adapt\'{e}e pour un livre de pr\'{e}paration au baccalaur\'{e}at.}
\smallskip

%----------------------------------------------------------
%  TABLEAU DES COMPOSANTES
%----------------------------------------------------------
\begin{center}
\renewcommand{\arraystretch}{1.3}
\begin{tabular}{|c|l|c|}
  \hline
  \rowcolor{navylight}
  \textbf{\color{navy}Exercice} & \multicolumn{1}{c|}{\textbf{\color{navy}Th\`{e}me}} & \textbf{\color{navy}Points} \\
  \hline
  1 & Suites num\'{e}riques & 4 \\
  \hline
  2 & Probabilit\'{e}s & 1,5 \\
  \hline
  3 & Probabilit\'{e}s & 2,5 \\
  \hline
  4 & \'{E}tude graphique & 4,75 \\
  \hline
  5 & Fonction et int\'{e}grale & 7,25 \\
  \hline
\end{tabular}
\end{center}

\bigskip

%==========================================================
%  EXERCICE 1 — SUITES NUMÉRIQUES  (4 pts)
%==========================================================
\noindent\textbf{Exercice 1 :}\\
Soit $(u_n)$ une suite num\'{e}rique d\'{e}finie par :
\[
  u_0 = 0 \qquad \text{et} \qquad
  u_{n+1} = \frac{4u_n + 4}{u_n + 4}
  \quad \text{pour tout } n \in \N
\]

\begin{enumerate}[leftmargin=1.5em, label=\textbf{\arabic*)}, itemsep=2pt, topsep=4pt]
  \item Calculer $u_1$ et $u_2$.
  \item Montrer par r\'{e}currence que $0 \leq u_n \leq 2$ pour tout $n \in \N$.
  \item
    \begin{enumerate}[label=\textbf{\alph*)}, itemsep=1pt, topsep=1pt]
      \item V\'{e}rifier que $u_{n+1} - 2 = \dfrac{2(u_n - 2)}{u_n + 4}$ pour tout $n \in \N$.
      \item Montrer que $u_{n+1} - u_n = \dfrac{(2 - u_n)(u_n + 2)}{u_n + 4}$,
            puis d\'{e}duire que $(u_n)$ est croissante.
      \item En d\'{e}duire que $(u_n)$ est convergente.
    \end{enumerate}
  \item On consid\`{e}re $(v_n)$ d\'{e}finie par $v_n = \dfrac{u_n - 2}{u_n + 2}$ pour tout $n \in \N$.
    \begin{enumerate}[label=\textbf{\alph*)}, itemsep=1pt, topsep=1pt]
      \item Montrer que $(v_n)$ est une suite g\'{e}om\'{e}trique de raison $\dfrac{1}{3}$.
      \item D\'{e}terminer $v_n$ en fonction de $n$, puis $u_n$ en fonction de $n$.
      \item Calculer $\displaystyle\lim_{n \to +\infty} u_n$.
      \item Calculer $S_n = v_0 + v_1 + \cdots + v_n$ en fonction de $n$.
    \end{enumerate}
\end{enumerate}

\medskip

%==========================================================
%  EXERCICE 2 — PROBABILITÉS  (1,5 pts)
%==========================================================
\noindent\textbf{Exercice 2 :}\\
Soit $X$ une variable al\'{e}atoire dont la loi de probabilit\'{e} est :

\medskip
\begin{center}
\renewcommand{\arraystretch}{1.4}
\begin{tabular}{|c|c|c|c|}
  \hline
  \rowcolor{navylight}
  \textbf{\color{navy}$k$} & $0$ & $1$ & $3$ \\
  \hline
  \textbf{\color{navy}$p(X=k)$} & $3a$ & $2a$ & $a$ \\
  \hline
\end{tabular}
\end{center}
\medskip

\begin{enumerate}[leftmargin=1.5em, label=\textbf{\arabic*.}, itemsep=2pt, topsep=4pt]
  \item D\'{e}terminer $a$.
  \item Calculer $E(X)$ et $V(X)$.
\end{enumerate}

\medskip

%==========================================================
%  EXERCICE 3 — PROBABILITÉS  (2,5 pts)
%==========================================================
\noindent\textbf{Exercice 3 :}\\
Une urne contient 5 boules noires num\'{e}rot\'{e}es $1;2;2;2;3$
et 4 boules blanches num\'{e}rot\'{e}es $1;1;2;3$.
On tire simultan\'{e}ment 3 boules. On pose :
\begin{itemize}[itemsep=0pt, topsep=2pt, leftmargin=2em]
  \item $A$ : \guillemotleft{} les 3 boules tir\'{e}es sont de m\^{e}me couleur \guillemotright{}
  \item $B$ : \guillemotleft{} les 3 boules tir\'{e}es portent le m\^{e}me num\'{e}ro \guillemotright{}
\end{itemize}

\begin{enumerate}[leftmargin=1.5em, label=\textbf{\arabic*.}, itemsep=2pt, topsep=4pt]
  \item Calculer $p(A)$, $p(B)$ puis montrer que $p(A \cap B) = \dfrac{1}{84}$.
  \item Les \'{e}v\'{e}nements $A$ et $B$ sont-ils ind\'{e}pendants ?
  \item Sachant que les boules tir\'{e}es sont de m\^{e}me couleur,
        calculer la probabilit\'{e} qu'elles portent le m\^{e}me num\'{e}ro.
\end{enumerate}

\medskip

%==========================================================
%  EXERCICE 4 — ÉTUDE GRAPHIQUE  (4,75 pts)
%==========================================================
\noindent\textbf{Exercice 4 :}\\
On consid\`{e}re la fonction $f$ d\'{e}finie par sa courbe repr\'{e}sentative $(\mathcal{C})$
dans un rep\`{e}re orthonorm\'{e}, avec une asymptote verticale $x = -3$
et une asymptote oblique $y = \dfrac{x+1}{2}$.\\
R\'{e}pondre graphiquement aux questions suivantes :

\begin{enumerate}[leftmargin=1.5em, label=\textbf{\arabic*.}, itemsep=2pt, topsep=4pt]
  \item D\'{e}terminer le domaine de d\'{e}finition de $f$.
  \item D\'{e}terminer le signe de $f$ sur son domaine.
  \item
    \begin{enumerate}[label=\textbf{\alph*.}, itemsep=1pt, topsep=1pt]
      \item D\'{e}terminer $\displaystyle\lim_{x \to +\infty} f(x)$
            et $\displaystyle\lim_{x \to -3^{+}} f(x)$.
      \item D\'{e}terminer les branches infinies de $(\mathcal{C})$.
    \end{enumerate}
\end{enumerate}

\newpage

%==========================================================
%  EXERCICE 5 — ÉTUDE DE FONCTION  (7,25 pts)
%==========================================================
\noindent\textbf{Exercice 5 :}\\
Soit $h$ la fonction d\'{e}finie sur $\R$ par $h(x) = (x-1)\,e^{x}$.
Soit $(\mathcal{C}_h)$ sa courbe repr\'{e}sentative dans un rep\`{e}re
orthonorm\'{e} $(O;\,\vec{i};\,\vec{j})$, avec $\|\vec{i}\| = 1\,\text{cm}$.

\begin{enumerate}[leftmargin=1.5em, label=\textbf{\arabic*.}, itemsep=2pt, topsep=4pt]
  \item Calculer $\displaystyle\lim_{x \to +\infty} h(x)$ et
        $\displaystyle\lim_{x \to +\infty} \dfrac{h(x)}{x}$,
        puis donner une interpr\'{e}tation g\'{e}om\'{e}trique.
  \item V\'{e}rifier que $h(x) = x\,e^{x}\!\left(1 - \dfrac{1}{x}\right)$
        pour tout $x \neq 0$, puis calculer $\displaystyle\lim_{x \to -\infty} h(x)$
        et donner une interpr\'{e}tation g\'{e}om\'{e}trique.
  \item
    \begin{enumerate}[label=\textbf{\alph*.}, itemsep=1pt, topsep=1pt]
      \item Montrer que $\forall x \in \R$ :
            $h'(x) = x\,e^{x}$.
      \item Dresser le tableau de variations de $h$ sur $\R$.
    \end{enumerate}
  \item D\'{e}terminer l'\'{e}quation de la tangente $(T)$ \`{a}
        $(\mathcal{C}_h)$ au point d'abscisse $0$.
  \item Montrer que la fonction $G : x \mapsto (x-2)\,e^{x}$
        est une primitive de $h$ sur $\R$.
  \item Calculer l'aire (en cm\textsuperscript{2}) du domaine plan d\'{e}limit\'{e}
        par $(\mathcal{C}_h)$, l'axe des abscisses et les droites
        $x = 1$ et $x = 3$.
  \item Montrer que $h(x) \geq -1$ pour tout $x \in \R$.
        En d\'{e}duire la position de $(\mathcal{C}_h)$ par rapport
        \`{a} la tangente $(T)$ sur $\R$.
\end{enumerate}

\begin{center}
  \color{charcoal}\rule{5cm}{0.4pt}\\[3pt]
  \textit{--- Fin du sujet ---}\\
  \rule{5cm}{0.4pt}
\end{center}


\end{document}